\documentclass[a4paper,14pt,russian]{extreport}
\usepackage{extsizes}
\usepackage{cmap} 
\usepackage[T2A]{fontenc}
\usepackage[utf8]{inputenc}
\usepackage[russian]{babel}
\usepackage{listings}
\usepackage{graphicx} 
\usepackage{amssymb,amsfonts,amsmath,amsthm} 
\usepackage{indentfirst} 
\usepackage[usenames,dvipsnames]{color}
\usepackage{makecell}
\usepackage{multirow}
\usepackage{ulem}
\usepackage{geometry}
\geometry{left=3cm}
\geometry{right=1.5cm}
\geometry{top=2.4cm}
\geometry{bottom=2.4cm}
\usepackage{fancyhdr}
\pagestyle{fancy}
\fancyhf{}
\fancyhead[R]{\thepage}
\fancyheadoffset{0mm}
\fancyfootoffset{0mm}
\setlength{\headheight}{17pt}
\renewcommand{\headrulewidth}{0pt}
\renewcommand{\footrulewidth}{0pt}
\fancypagestyle{plain}{
	\fancyhf{}
	\rhead{\thepage}}
\setcounter{page}{1}
\usepackage{float}
\usepackage{titlesec}
\titleformat{\chapter}[display]
{\filcenter}
{\MakeUppercase{\chaptertitlename} \thechapter}
{8pt}
{\bfseries}{}
\titleformat{\section}
{\normalsize\bfseries}
{\thesection}
{1em}{}
\titleformat{\subsection}
{\normalsize\bfseries}
{\thesubsection}
{1em}{}
\titlespacing*{\chapter}{0pt}{-30pt}{8pt}
\titlespacing*{\section}{\parindent}{*4}{*4}
\titlespacing*{\subsection}{\parindent}{*4}{*4}
\usepackage{tocloft}
\renewcommand{\cfttoctitlefont}{\hspace{0.38\textwidth}
	\bfseries\MakeUppercase}
\renewcommand{\cftbeforetoctitleskip}{-1em}
\renewcommand{\cftaftertoctitle}{\mbox{}\hfill \\
	\mbox{}\hfill{\footnotesize Стр.}\vspace{-2.5em}}
\renewcommand{\cftchapfont}{\normalsize\bfseries
	\MakeUppercase{\chaptername}}
\renewcommand{\cftsecfont}{\hspace{31pt}}
\renewcommand{\cftsubsecfont}{\hspace{11pt}}
\renewcommand{\cftbeforechapskip}{1em}
\renewcommand{\cftparskip}{-1mm}
\renewcommand{\cftdotsep}{1}
\setcounter{tocdepth}{2}
\newcommand{\empline}{\mbox{}\newline}
\newcommand{\likechapterheading}[1]{
	\begin{center}
		\textbf{\MakeUppercase{#1}}
	\end{center}
	\empline}
\makeatletter
\renewcommand{\@dotsep}{2}
\newcommand{\l@likechapter}[2]
{{\bfseries\@dottedtocline{0}{0pt}{0pt}{#1}{#2}}}
\makeatother
\newcommand{\likechapter}[1]{
	\likechapterheading{#1}
	\addcontentsline{toc}{likechapter}{\MakeUppercase{#1}}}
\usepackage[tableposition=top]{caption}
\usepackage{subcaption}
\DeclareCaptionLabelFormat{gostfigure}{Рисунок #2}
\DeclareCaptionLabelFormat{gosttable}{Таблица #2}
\DeclareCaptionLabelSeparator{gost}{~---~}
\captionsetup{labelsep=gost}
\captionsetup[figure]{labelformat=gostfigure}
\captionsetup[table]{labelformat=gosttable}
\renewcommand{\thesubfigure}{\asbuk{subfigure}}
\usepackage[square,numbers,sort&compress]{natbib}
\renewcommand{\bibnumfmt}[1]{#1.\hfill}
\renewcommand{\bibsection}{\likechapter{Список литературы}}
\setlength{\bibsep}{0pt}
\usepackage{changepage}
\linespread{1.3}
%\renewcommand{\rmdefault}{ftm}
\renewcommand{\_}{\kern-1.5pt\textunderscore\kern-1.5pt}
\frenchspacing
\begin{document}
	\begin{minipage}[b]{0.2\textwidth}
		\includegraphics[width=1\linewidth]{pics/image1.png}
	\end{minipage}
	\begin{minipage}[b]{0.8\textwidth}
		\begin{center}
			\begin{small}
				\textbf{Институциональные реквизиты удалены}
			\end{small}
		\end{center}
	\end{minipage}
	
	\begin{adjustwidth}{0.59in}{0.0in}
		\textbf{Подразделение удалено\par}
	\end{adjustwidth}
	
	\begin{adjustwidth}{2.56in}{0.0in}
		\begin{flushright}
			Допустить к защите
		\end{flushright}\par
	\end{adjustwidth}
	
	\begin{adjustwidth}{2.56in}{0.0in}
		\begin{flushright}
			зав. кафедрой
		\end{flushright}\par
	\end{adjustwidth}
	
	\begin{adjustwidth}{2.56in}{0.0in}
		\begin{center}
			\begin{small}
				\ \ \ \ \ \ \ \ \ \ \ \  \uline{\ \ \ \ \ \ \ \ \ \ \ \ \ \ \ \ \ \ \ \ \ \ \ \ \ \ \ \ \ \ \ \ \ \ \ \ \  } / \uline{ Яшина М.В. } /
			\end{small}
		\end{center}\par
	\end{adjustwidth}
	
	\begin{adjustwidth}{2.56in}{0.0in}
		\begin{center}
			\ \ \ \ \ \ \ \ \ \ \ \ \ \ \ \ \ \  \  « \uline{\ \ \ \ \  }»\uline{\ \ \ \ \ \ \ \ \ \ \ \ \ \ \ \ \ \ \ \ \ \ \ \ \ \ \ \ \ \  }20\uline{\ \ \ \  }г.
		\end{center}\par
	\end{adjustwidth}
	
	\begin{center}
		\begin{large}
			Халилов Джабраиль Эльнурович\par
		\end{large}
	\end{center}
	
	\begin{center}
		\textbf{ВЫПУСКНАЯ КВАЛИФИКАЦИОННАЯ РАБОТА}
	\end{center}\par
	
	\begin{flushleft}
		на тему: "Сравнение различных вариационных принципов построения расчетных сеток" \ \ \ \ \ \ \ \ \ \ \ \ \ \ \ \ \ \ \ \ \ \ \ \ \ \ \ \ \ \ \ \ \ \ \ \ \ \ \ \ \ \ \ \ \ \ \ \ \ \ \ \ \ \ \ \ \ \ \ \ \ \ \ \ \ \ \ \ \ \ \ \ \ \ \ \ \ \ \ \ \ \ \  
		\\
		
		шифр и направление подготовки \par
		01.03.04 Прикладная математика \ \ \ \ \ \ \ \ \ \ \ \ \par
		\setlength{\parskip}{0.0pt}
		\textcolor{white}{ффф}\\
		направленность (профиль) \  "Прикладная математика"  \par
		
		группа\ \ \  4бПМ  \par
		
	\end{flushleft}
	
	\begin{flushleft}
		Автор ВКР\ \ \ \ \ \ \ \ \ \ \ \ \ \ \ \ \ \ \ \ \  \uline{\ \ \ \ \  Халилов Д.Э. \ \ \ \ \ \ \ \ \  }/\uline{\ \ \ \ \ \ \ \ \ \ \ \ \ \ \ \ \ \ \ \ \ \ \ \ \  }/\par
		
		{\fontsize{10pt}{12.0pt}\selectfont \ \ \ \ \ \ \ \ \ \ \ \ \ \ \ \ \ \ \ \ \ \ \ \ \ \ \ \ \ \ \ \ \ \ \ \ \ \ \ \ \ \ \ \ \ \ \ \ \ \ \ \ \ \ \ \ \ \ \ \ \ \ \ \ \ \ \ \ \ \ \ \ \ \ \ \ \ \  (Ф.И.О., подпись, дата)\par}
		
		Руководитель ВКР\ \ \ \ \ \ \ \ \ \uline{\ \ \ \ \ \ Тишкин В.Ф. \ \ \ \ \ \ \ \  }/\uline{\ \ \ \ \ \ \ \ \ \ \ \ \ \ \ \ \ \ \ \ \ \ \ \ \  }/\par
		
		{\fontsize{10pt}{12.0pt}\selectfont \ \ \ \ \ \ \ \ \ \ \ \ \ \ \ \ \ \ \ \ \ \ \ \ \ \ \ \ \ \ \ \ \ \ \ \ \ \ \ \ \ \ \ \ \ \ \ \ \ \ \ \ \ \ \ \ \ \ \ \ \ \ \ \ \ \ \ \ \ \ \ \ \ \ \ \ \ \  (Ф.И.О., подпись, дата)\par}
		
		Консультант (ы) ВКР\  \ \ \ \ \  \uline{\ \ \ \ \ Самохвалова Ж.П. \ \ \ \ \ \ \ \  }/\uline{\ \ \ \ \ \ \ \ \ \ \ \ \ \ \ \ \ \ \ \ \ \ \ \ \ \  }/\par
		
		{\fontsize{10pt}{12.0pt}\selectfont \ \ \ \ \ \ \ \ \ \ \ \ \ \ \ \ \ \ \ \ \ \ \ \ \ \ \ \ \ \ \ \ \ \ \ \ \ \ \ \ \ \ \ \ \ \ \ \ \ \ \ \ \ \ \ \ \ \ \ \ \ \ \ \ \ \ \ \ \ \ \ \ \ \ \ \ \ \  (Ф.И.О., подпись, дата)\par}
		
		
	\end{flushleft}
	\begin{small}
		\begin{center}
			2024
		\end{center}
	\end{small}
	
	\thispagestyle{empty}
	
	\newpage
	\begin{center}
		Институциональные реквизиты удалены
	\end{center}
	Подразделение удалено\\
	Шифр и направление подготовки\\
	\ \textcolor{white}{п}\uline{\ \ \ \ \ \ \ \ \ \ \ \ \ \ \ \ \ \ 01.03.04 Прикладная математика \ \ \ \ \ \ \ \ \ \ \ \ \ \ \ \ \ \ \ \ \ \ \ \ \ \ \ \ \ \ \ \ \ \ \ \ \ \ \ \ \ \ \ \ \ \ \ \ \ \ \  }\\
	Направленность (профиль)/\ \textcolor{white}{п}\uline{ \ \ \ "Прикладная математика" \ \ \ \ \ \ \ }\\
	
	\begin{center}
		ЗАДАНИЕ \\ НА ВЫПУСКНУЮ КВАЛИФИКАЦИОННУЮ РАБОТУ\\
		Халилов Джабраиль Эльнурович\\
		(Фамилия, Имя, Отчество обучающегося)
	\end{center}
	
	\begin{flushleft}
		Тема ВКР: \uline{Сравнение различных вариационных принципов построения
расчетных сеток}\ \ \ \ \ \ \ \ \ \ \ \ \ \ \ \ \ \ \ \ \ \ \ \ \ \ \ \ \ \ \ \ \ \ \ \ \ \ \ \ \ \ \ \ \ \ \ \ \ \ \ \ \ \ \ \ \ \ \ \ \ \ \ \ \ \ \ \ \ \ \ \ \ \ \ \ \ \ \ \ \ \ \ \ \ \  
		\\
		институциональные реквизиты утверждения удалены \\
		1. Исходные данные по ВКР:\\
		Характер работы --- научно-исследовательская работа.\\

\begin{flushleft}
\thispagestyle{empty}
2. Обоснование темы ВКР и перечень, подлежащих разработке вопросов:\\
\uline{Построение\ расчётных\ сеток\ —\ ключевой\ элемент\ численного\ решения\ уравнений\ с\ частными\ производными.\ Программная\ реализация\ на\ Python\ с\ использованием\ стандартных\ библиотек.\ Требуется\ анализ\ и\ сравнение\ вариационных\ методов,\ оценка\ точности\ и\ устойчивости\ сеток.\ Результаты\ анализа\ документировать\ и\ сохранять\ в\ txt\ для\ дальнейшего\ построения\ графиков\ и\ анализа.}\\
\end{flushleft}
		
		\thispagestyle{empty}
\begin{table}[ht]
    \centering
    \begin{tabular}{|c|c|c|c|c|c|}
        \hline
        \small{№ п/п} & \small{Наименование}  & \small{Ф.И.О. должность,} & \small{Срок}  &\multicolumn{2}{|c|}{\small{подпись дата}} \\ 
        \cline{5-6}
        
        & \small{этапа} & \small{уч. степень, звание} & \small{выполнения} &\small{Задание} & \small{Задание}\\ 
        & \small{работы}  & \small{руководителя/} & \small{этапа,} & \small{выдал} & \small{принял} \\ 
        &\small{(раздела)} & \small{консультанта} & \small{раздела} & & \\ 
        & & \small{(этапа работы раздела)}& & & \\ \hline
        
        
        \small{1}  & \small{Подготовка}  & \small{Тишкин В.Ф., }& \small{01.02.2024} & & \\ 
        & \small{теоретических} &  \small{чл.-корр. РАН,} & & & \\
        & \small{сведений и } & \small{ д.ф.-м.н., проф.} &  & & \\  
        & \small{изучение материалов} &  &  & & \\
        \hline  
        
        \small{2}  & \small{Создание и отладка} & \small{Тишкин В.Ф., } & \small{01.03.2024} & & \\ 
        &  \small{приложения} & \small{чл.-корр. РАН,} &  & & \\
        & & \small{д.ф.-м.н., проф.} & & & \\ \hline

        \small{3}  & \small{Написание экономи-} & \small{Самохвалова Ж.П.,}  & \small{01.04.2024}  & & \\ 
        & \small{ческого обоснования} &\small{ ст.пр. } &  & & \\ \hline  

        \small{4}  & \small{Написание} & \small{Тишкин В.Ф., } & \small{01.05.2024} & & \\ 
        & \small{текста диплома} & \small{чл.-корр. РАН,} & & & \\
        & & \small{д.ф.-м.н., проф.} & & & \\ \hline
        
        \small{5}  & \small{Оформление} & \small{Тишкин В.Ф., }&\small{ 15.06.2024} & & \\
        & \small{текста ВКР} &\small{чл.-корр. РАН,} & & & \\
        & & \small{д.ф.-м.н., проф.}& & & \\
        \hline
    \end{tabular}
\end{table}

		\begin{small}
			Руководитель ВКР\ \ \ \ \ \ \ \ \ \ \ \ \ \ \ \ \ \ \ \uline{\ \ \ \ \ \ \ \ \ \ \ \ Тишкин В.Ф. \ \ \ \ \ \ \ \ \ \ \ \  }/\uline{\ \ \ \ \ \ \ \ \ \ \ \ \ \ \ \ \ \ \ \ \ \ \ \  }/\par
			
			{\fontsize{10pt}{12.0pt}\selectfont \ \ \ \ \ \ \ \ \ \ \ \ \ \ \ \ \ \ \ \ \ \ \ \ \ \ \ \ \ \ \ \ \ \ \ \ \ \ \ \ \ \ \ \ \ \ \ \ \ \ \ \ \ \ \ \ \ \ \ \ \ \ \ \ \ \ \ \ \ \ \ \ \ \ \ \ \ \  (Ф.И.О., подпись, дата)\par}
			Задание принял\\
			к исполнению\ \ \ \ \  \ \ \ \ \ \ \ \ \ \ \ \ \ \ \ \ \ \ \ \ \  \uline{\ \ \ \  \ \ \ \ \  \ Халилов Д. Э. \ \ \ \ \ \ \ \ \ \ \ \ \ \ \ \  }/\uline{\ \ \ \ \ \ \ \ \ \ \ \ \ \ \ \ \ \ \ \ \ \ \ \  }/\par
			
			{\fontsize{10pt}{12.0pt}\selectfont \ \ \ \ \ \ \ \ \ \ \ \ \ \ \ \ \ \ \ \ \ \ \ \ \ \ \ \ \ \ \ \ \ \ \ \ \ \ \ \ \ \ \ \ \ \ \ \ \ \ \ \ \ \ \ \ \ \ \ \ \ \ \ \ \ \ \ \ \ \ \ \ \ \ \ \ \ \  (Ф.И.О., подпись, дата)\par}
			
		\end{small}
		
	\end{flushleft}
	\newpage
	
	\thispagestyle{empty}
	\begin{center}
		ОТЗЫВ РУКОВОДИТЕЛЯ НА ВЫПУСКНУЮ КВАЛИФИКАЦИОННУЮ РАБОТУ
	\end{center}
	Обучающейся
	\ \uline{\ \ \ Халилов Джабраиль Эльнурович \ \ \ \ \ \ \ \ \ \ \ \ \ \ \ \ \ \ \ \ \ \ \ \ \ \ \ \ \ \ \ \ \ \ \ \ \ \ \ \ \  }\\
	Шифр и направление подготовки/специальность\\
	\ \uline{\ \ \ \ \ \ \ \ \ \ \ \ \ \ \ \ \ \ 01.03.04 Прикладная математика \ \ \ \ \ \ \ \ \ \ \ \ \ \ \ \ \ \ \ \ \ \ \ \ \ \ \ \ \ \ \ \ \ \ \ \ \ \ \ \ \ \ \ \ \ \ \ \ \ \ \  }\\
	Направленность (профиль)/ специализация \ \textcolor{white}{п}\uline{ \ \ \ Прикладная математика \ \ \ \ \ \ \ }\\
	Группа \ \uline{ \ \ \ 4бПМ \ \ \ \ \ \ \ }\\
	Руководитель ВКР \ \textcolor{white}{п}\uline{ \ \ \ Тишкин В.Ф., чл.-корр. РАН, д.ф-м.н., проф. \ \ \ \ \ \ \ }\\
	{\fontsize{10pt}{12.0pt}\selectfont \ \ \ \ \ \ \ \ \ \ \ \ \ \ \ \ \ \ \ \ \ \ \ \ \ \ \ \ \ \ \ \ \ \ \ \ \ \ \ \  (Ф.И.О., ученая степень и (или) ученое звание)\par}
	Тема: \uline{Сравнение различных вариационных принципов построения}\\
    \uline{расчетных сеток} \ \ \ \ \ \ \ \ \ \ \ \ \ \ \ \ \ \ \ \ \ \ \ \ \ \ \ \ \ \ \ \ \ \ \ \ \ \ \ \ \ \ \ \ \ \ \ \ \ \ \ \ \ \ \ \ \ \ \ \ \ \ \ \ \ \ \ \  
	\\
	Содержание отзыва: \\
 \uline{Выпускная квалификационная работа Халилова Д.Э. относится к актуальной тематике~--- построению расчётных сеток для численного моделирования.\\ 
    В работе проведена аналитическая работа и рассмотрены три различных метода построения расчётных сеток.
    Описаны основные принципы каждого метода, а также их преимущества и недостатки. 
    Проведено сравнение методов на основе двух ключевых критериев: коэффициента ортогональности и критерия равномерности. Представлены результаты сравнительного анализа, которые демонстрируют эффективность каждого метода.
    Приведены графики и их анализ, подтверждающие выводы исследования.
    Считаю, что работа Халилова Д.Э. удовлетворяет требованиям,
    предъявляемым к выпускным квалификационным работам, и заслуживает высокой оценки.}

	\uline{}
	
	\begin{small}
		
		\thispagestyle{empty}
		\ \ \ \ \ \ \uline{Тишкин В.Ф.} \ \ \ \ \ \ \ \ \ \ \ \ \ \ \ \ \ \ \ \ \ \uline{\ \ \ \ \ \ \ \ \ \ \ \ \ \ \ \ \ } \ \ \ \ \ \ \ \ \ \ \ \ \ \ \ \uline{\ \ \ \ \ \ \ \ \ \ \ \ \ \ \ \ \ \ \ \ \ \ \ \ }\par
		{\fontsize{10pt}{12.0pt}\selectfont  (Ф.И.О. руководителя ВКР) \ \ \ \ \ \ \ \ \ \ \ \ \ \ \ \ \ \ \ \ \ \ \ \ (дата) \ \ \ \ \ \ \ \ \ \ \ \ \ \ \ \ \ \ \ \ \ \ \ \ \ \ \ \  \ \ \ (Подпись) }\\
		С отзывом ознакомлен:\par
		\ \ \ \ \ \ \uline{Халилов Д.Э.} \ \ \ \ \ \ \ \ \ \ \ \ \ \ \ \ \ \ \ \ \ \uline{\ \ \ \ \ \ \ \ \ \ \ \ \ \ \ \ \ } \ \ \ \ \ \ \ \ \ \ \ \ \ \ \ \uline{\ \ \ \ \ \ \ \ \ \ \ \ \ \ \ \ \ \ \ \ \ \ \ \ }\par
		{\fontsize{10pt}{12.0pt}\selectfont \textcolor{white}{п}  (Ф.И.О. обучающегося) \ \ \ \ \ \ \ \ \ \ \ \ \ \ \ \ \ \ \ \ \ \ \ \ \ \ \ \ (дата) \ \ \ \ \ \ \ \ \ \ \ \ \ \ \ \ \ \ \ \ \ \ \ \ \ \ \ \ \ \ \ \ (Подпись) }\\
	\end{small}
	\newpage
	\begin{center}
		РЕФЕРАТ
  	\end{center} 
   
В процессе выполнения дипломной работы был проведен анализ и сопоставление различных принципов построения расчетных сеток. Расчетные сетки играют ключевую роль в численном моделировании, где точность и эффективность численных методов зависят от качества сетки. В работе были изучены основные методы создания структурированных сеток, включая принципы геометрической адаптации и оптимизации элементов сетки.

В рамках дипломной работы используются теоретические знания для разработки и оценки эффективности сеток в сложных геометрических областях, что способствует достижению высокой точности приближения решений прикладных задач. Особое внимание уделяется сравнению методов построения сеток на основе вариационных принципов, что включает разработку алгоритмов для автоматического создания сеток, анализ их структурных характеристик и оценку влияния на качество решений.

Исследование проводится при использовании передовых программных инструментов и языков программирования высокого уровня, что обеспечивает точность расчетов и удобство анализа результатов. Полученные данные представляют значительный вклад в развитие методов математического моделирования и могут быть использованы для будущих исследований в данной области.

Целью данной работы является не только теоретическое сопоставление различных подходов к созданию сеток, но также практическая разработка универсальных решений, способных подстраиваться под специфику конкретных прикладных задач.

В списке использованной литературы приводятся названия всех использованных 
источников в количестве 14 штук.
Пояснительная записка ВКР содержит 72 страницы, 11 рисунков, 3 приложения. 
	
	
	\newpage
	\renewcommand\contentsname{\centerline{СОДЕРЖАНИЕ}}
	\thispagestyle{empty}
	\tableofcontents
	\thispagestyle{empty}
	
\chapter{ВВЕДЕНИЕ}
    При решении многих прикладных задач численными методами важное место занимают расчетные сетки. Стадия построения такого плана сеток играет большую роль в общем ходе (математического) моделирования. В областях с сложной геометрией создание расчётной сетки является задачей, требующей специальных программных навыков, и качество получаемого приближенного решения во многом зависит от удачного выбора сетки. Построение расчётных сеток требует учёта различных факторов, таких как форма изучаемой области, необходимая точность решения и особенности численных методов. Важно подобрать размеры ячеек сетки таким образом, чтобы достигнуть оптимального баланса между точностью и вычислительной сложностью. Существуют специальные методы построения сеток, которые эффективно учитывают геометрические особенности для оптимального использования вычислительных ресурсов. Некоторые задачи требуют применения неструктурированных сеток, которые лучше адаптируются к форме области и распределению физических свойств, однако использование таких сеток требует более сложных численных алгоритмов и больших вычислительных мощностей. Таким образом, построение расчетных сеток представляет важный этап в численном моделировании, который учитывает различные факторы для достижения наилучших результатов. В данной работе для решения задач будет использоваться Python 3 со следующими библиотеками: NumPy и Matplotlib. NumPy - это мощная библиотека для математических вычислений, включая линейную алгебру; Matplotlib предоставляет собой библиотеку для визуализации данных. 
\par Далее мы продолжим рассматривать факторы, которые повляют на качество расчетных сеток и их функцию в различных сферах инженерии (проектирования) и прикладной математики.

	
\section{Постановка задачи}
\par В рамках данного исследования рассматривается генерация структурированных двумерных сеток в изогнутых четырехугольниках. Представляют они из себя связные области с границей, которые состоят из фрагментов плавной линии. Эти четырехугольные фигуры имеют четыре определенные точки на границе. Данные точки не всегда совпадают с точками изгиба.

\par Сначала рассмотрим метод, при котором узлы на границе заданы "до", а сетка создается с использованием нитей, что могут гнуться и растянуты таким образом, чтобы их краевые узлы совпадали заранее конкретными узлам четырехугольника. Данный процесс изгибания решетки приводит к тому, что формируется нужная нам сетка. В состоянии равновесия потенциальная энергия системы становится минимальной. Происходит это в результате уравнений для внутренних узлов сетки.


\par Процесс уменьшения потенциальной энергии приводит к набору уравнений, решив, которые мы сможем найти расположение узлов внутри и закончить построение сеток. Данные уравнения - это линейные системы, методы решения которых подробно изучены. При работе с не выпуклыми областями важно учитывать различную жесткость упругих элементов, чтобы избежать выхода узлов за пределы области и возникновения самопересечений сетки.

\par Подводя итог, анализ изменений метода, таких как изменение параметров упругости элементов, позволит нам приспособить сетку к особенностям конкретной задачи, что обеспечит построение сложных геометрических фигур и форм.


\section{Методы и подходы к построению сеток}
\par Первый метод, метод включает в себя разбор критериев минимизации функционалов, адаптацию сеток в условиях некоторых физических и геометрических ограничений. Предложенный подход поможет продуктивно применить структурированные сетки в плоских криволинейных областях, что позволит расширить возможности их применения в различных научных задачах и инженерных проблем. Для краткости и комфорта будем называть этот метод - алгоритмический. 
\par Второй метод, известен как метод Винслоу, назван он так по фамилии своего создателя, метод представляет из себя модификацию стандартных подходов, что продуктивно отобразится при использовании с различными трехмерными пространствами. Данный метод позволяет преодолевать ограничения предыдущего подхода, что дает возможность создания более комплексных и достаточно, выверенных структур для различных приложений. 
 Метод включает криволинейные четырехугольники с конкретно заданным расположением узлов на границах данной фигуры это позволяет формировать стационарные температуры полей, что дает возможность использовать метод в задачах, где требуется четкое соблюдение граничных условий. Выводы программ, извлекаемые посредством данного метода, опираются на минимизации функционала Дирихле и это приводит к образованию изотерм, которые показывают формирование. Данные изотермы мы используем для создания точных и адаптированных сеток, подобного рода сетки способны показать сложные геометрические формы.
\par Третий способ, предложенный моим научным руководителем, который я предлагаю назвать “Метод адаптивного натяжения” для удобства в данной работе, представляет собой подход к созданию сеток, основанный на минимизации функционала, связанного с матрицей Якоби и ее собственными значениями. Основная идея этого метода заключается в том, что для каждой точки сетки вычисляется матрица Якоби, определяются ее собственные значения и минимизируется функционал, содержащий два компонента: сумму квадратов разностей собственных значений матрицы Якоби, деленную на определитель матрицы, и квадрат разности определителя матрицы Якоби и единицы. Итерационный процесс минимизации функционала позволяет обновлять координаты точек сетки таким образом, чтобы уменьшить значение функционала. Это способствует формированию качественных структур сетки для сложных областей. Кроме того, метод учитывает граничные условия для создания строго соответствующих заданной геометрии области. Применение данного метода оказывается эффективным в различных инженерных и научных задачах, требующих точного выполнения граничных условий и учета геометрических особенностей.

\chapter{Методы построения}
\section{Использование отображений для построения структурированных сеток}

\par Рассмотрим процесс создания двухмерных структурированных сеток в области, известной как криволинейный четырехугольник. Под криволинейным четырехугольником понимается область с непрерывной границей, на которой определены четыре точки, называемые углами криволинейного четырехугольника. Эти точки могут не совпадать с изломами границы. На рисунке 2.1 можно увидеть примеры таких криволинейных четырехугольников.
\begin{figure}[h!]
\centering
\includegraphics[width=0.8\textwidth]{pics/ris2.1.png} 
\caption{Примеры криволинейных четырехугольников}
\label{fig:curved_quadrilaterals}
\end{figure}
	
\par Давайте начнем с одного из основных способов построения сетки внутри криволинейного четырехугольника, если мы знаем расположение его граничных узлов. Этот метод использует сетку из упругих нитей, соединенных в узловых точках. При растяжении такой сетки мы перемещаем ее граничные узлы так, чтобы они соответствовали узлам криволинейного четырехугольника, как показано на рисунке 2.2. В процессе растягивания сетка изменяется формой, образуя нужную структуру сетки.

\begin{figure}[H]
  \centering
  \includegraphics[width=0.8\textwidth]{pics/ris2.2.png}
  \caption{Пример деформированной сетки в криволинейном четырехугольнике}
  \label{fig:curved_quadrilaterals}
\end{figure}

    \par Данный четырехугольник со связкой упругих нитей. Решетка, которая натянута на четырехугольник. Теперь нужно определить местоположение узлов деформированной решетки. Для этого применим принцип, согласно которому потенциальная энергия системы в состоянии равновесия минимальна. Как известно из физики, потенциальная энергия растянутого упругого элемента определяется следующим образом: 
    
\begin{equation}
\Pi = k \frac{(l - l_0)^2}{2}
\end{equation}


где \( k \) — коэффициент упругости, \( l \) — текущая длина элемента, а \( l_0 \) — его исходная длина.
Будем предполагать, что 
\[
l_0 \ll l.
\]
Тогда все можно упростить и получить
\begin{equation}
\Pi=\mathrm{k} \frac{l^2}{2}.
\end{equation}
Из этого факта также следует, что все упругие элементы в деформируемой сетке растягиваются, образуя прямолинейные отрезки. Это важно, потому что каждый элемент сетки, представляющий собой упругую нить, стремится к минимизации своей потенциальной энергии. Поэтому в процессе деформации сетки каждый из этих элементов выстраивается в прямую линию между узлами, к которым он прикреплен.

Этот принцип позволяет нам утверждать, что несмотря на криволинейную форму четырехугольника, элементы деформированной сетки всегда будут прямыми. Это свойство значительно облегчает процесс моделирования и расчета, поскольку можно оперировать линейными элементами даже в условиях сложных геометрических форм.

Более того, растягивание сетки в соответствии с указанным принципом минимизации энергии приводит к тому, что сетка адаптируется к границам криволинейного четырехугольника и при этом сохраняется упорядоченной и структурированной формой. Это способствует равномерному распределению узлов и точному представлению области.
\begin{figure}[H]
  \centering
  \includegraphics[width=0.8\textwidth]{pics/ris2.3.png}
  \caption{Индексация узлов}
  \label{fig:curved_quadrilaterals}
\end{figure}
Индексация узлов и элементов деформированной решетки осуществляется таким образом: узлы, принадлежащие к сетке, нумеруются двумя индексами i и j, где i принимает значения от 0 до N, а j - от 0 до M. Рёбра сетки (упругие элементы деформированной решетки) обозначаются полуцелыми индексами i+1/2 и j, и i и j+1/2.

Данный способ индексации дает нам возможность точно определять позицию каждого узла и каждого ребра в деформированной сетке, что упростит следующие вычисления или моделирования.

\begin{equation}
\Pi_{i j+1 / 2}=\mathrm{k} \frac{\left(x_{i j+1}-x_{i j}\right)^2+\left(y_{i j+1}-y_{i j}\right)^2}{2}
\end{equation}

Полная упругая энергия системы представляет собой сложение энергий всех ребер по отдельности. Так как мы установили фиксацию граничных узлов, вся энергия деформации зависит исключительно от координат внутренних узлов сетки.

В итоге, если мы будем рассматривать сетку как совокупность упругих элементов, то вся потенциальная энергия системы будет определяться расположением внутренние узлы, а это позволит посчитать полную упругую энергию как функцию координат внутренних узлов. Данный факт является ключевым моментом при минимизации энергии для нахождения оптимальной конфигурации сетки.
\begin{align*}
\Pi_{\text{полн}} = & \frac{k}{2} \sum_{i=0}^{N-1} \sum_{j=0}^M \left[(x_{i+1,j} - x_{i,j})^2 + (y_{i+1,j} - y_{i,j})^2\right] + \\
& \frac{k}{2} \sum_{j=0}^{M-1} \sum_{i=0}^N \left[(x_{i,j+1} - x_{i,j})^2 + (y_{i,j+1} - y_{i,j})^2\right]
\end{align*}

Для достижения минимальной потенциальной энергии необходимо, чтобы частные производные функции по её аргументам были равны нулю. Это означает, что внутренние узлы сетки должны быть расположены таким образом, чтобы их изменение координат не приводило к изменению общей энергии системы. Если все частные производные равны нулю, то система находится в равновесии и потенциальная энергия достигает своего минимума.
\begin{equation}
    \frac{\partial \Pi_{\text{полн}}}{\partial x_{ij}} = 0, \quad i=1, \ldots, N-1
\end{equation}
\begin{equation}
    \frac{\partial \Pi_{\text{полн}}}{\partial y_{ij}} = 0, \quad j=1, \ldots, M-1
\end{equation}

Очевидно, что эти производные выглядят следующим образом:
\begin{equation}
    (x_{i+1,j} - 2x_{ij} + x_{i-1,j}) + (x_{i,j+1} - 2x_{ij} + x_{i,j-1}) = 0
\end{equation}
\begin{equation}
    (y_{i+1,j} - 2y_{ij} + y_{i-1,j}) + (y_{i,j+1} - 2y_{ij} + y_{i,j-1}) = 0 \quad
\end{equation}

Формулы представляют собой набор уравнений, которые, решив их, позволяют определить расположение внутренних точек и завершить создание сетки. Это линейные уравнения, для которых существуют эффективные методы решения.

Эту систему уравнений можно использовать для формирования сеток внутри выпуклых областей. Примеры таких сеток можно увидеть на иллюстрациях 2.4, 2.5 и 2.6. Однако при работе с невыпуклой областью координаты точек, полученные после решения уравнений, могут выходить за её границы. Для решения этой проблемы необходимо соответствующим образом корректировать коэффициенты жесткости “горизонтальных” и “вертикальных” элементов.

\begin{figure}[H]
  \centering
  \includegraphics[width=0.8\textwidth]{pics/ris2.5.png}
  \caption{Малый вертикальный коэффициент натяжения}
  \label{fig:curved_quadrilaterals}
\end{figure}

\begin{figure}[H]
  \centering
  \includegraphics[width=0.8\textwidth]{pics/ris2.4.png}
  \caption{Средний вертикальный коэффициент натяжения}
  \label{fig:curved_quadrilaterals}
\end{figure}

\begin{figure}[H]
  \centering
  \includegraphics[width=0.8\textwidth]{pics/ris2.6.png}
  \caption{Большой вертикальный коэффициент натяжения}
  \label{fig:curved_quadrilaterals}
\end{figure}




        
        \subsection{Невыпуклая область}
	\par В случае, когда невыпуклые области присутствуют, возникает проблема: упругие элементы одного типа могут быть недостаточно сильными, чтобы сопротивляться сжатию элементов другого типа и удерживать их внутри заданной области. Эту сложность можно решить путем усиления жесткости элементов данного семейства. Правильный выбор коэффициентов жесткости поможет избежать выхода узлов за пределы области и предотвратить пересечения сетки, как показано на рисунке 2.6.

    \par До этого мы использовали тот факт, что коэффициенты жесткости всех упругих элементов были одинаковые, но теперь мы принимаем во внимание, что жесткость горизонтальных и вертикальных элементов будет различна, причём соответствующие коэффициенты обозначены как $k_{\text{г}}$ и $k_{\text{в}}$ соответственно. Формулы для потенциальной энергии и уравнения для нахождения координат узлов также будут переписаны с учетом различий в коэффициентах жесткости:
    
    \begin{align}
    \Pi_{\text{полн}} &= \frac{k_{\Gamma}}{2} \sum_{i=0}^{N-1} \sum_{j=0}^{M} \left[\left(x_{i+1,j} - x_{ij}\right)^2 + \left(y_{i+1,j} - y_{ij}\right)^2\right] \notag \\
    &\quad + \frac{k_{B}}{2} \sum_{j=0}^{M-1} \sum_{i=0}^{N} \left[\left(x_{i,j+1} - x_{ij}\right)^2 + \left(y_{i,j+1} - y_{ij}\right)^2\right]
    \end{align}
    
    \begin{equation}
    k_{\Gamma} \left(x_{i+1,j} - 2x_{ij} + x_{i-1,j}\right) + k_{B} \left(x_{i,j+1} - 2x_{ij} + x_{i,j-1}\right) = 0
    \end{equation}
    

    \par Заметим, что можно отказаться от требования о том, чтобы узлы сетки на границе совпадали с заранее заданными точками. В таком случае нам придется добавить уравнения для граничных узлов (исключая четыре угла криволинейного четырехугольника, которые остаются неизменными). Данные уравнения отличаются от уравнений для внутренних узлов, так как граничные узлы не могут смещаться произвольным образом, а должны быть на границе, то есть координаты такого узла $x_{i_0} y_{i_0}$ становятся функционально связаны, а если считать, что граница, на которую попадают узлы $x_{i_0} y_{i_0}$ задается как неявная функция $F(x,y) = 0$, то соответствующее уравнение имеет вид:

$$
\begin{aligned}
& \mathrm{k}_{\Gamma}\left[\left.\left(x_{i+10}-2 x_{i 0}+x_{i-10}\right) \frac{\partial F}{\partial y}\right|_{x_{i 0}, y_{i 0}}-\left.\left(y_{i+10}-2 y_{i 0}+y_{i-10}\right) \frac{\partial F}{\partial x}\right|_{x_{i 0}, y_{i 0}}\right]+ \\
& +\mathrm{k}_{\mathrm{B}}\left[\left.\left(x_{i 1}-x_{00}\right) \frac{\partial F}{\partial y}\right|_{x_{i 0}, y_{i 0}}-\left.\left(y_{i 1}-y_{00}\right) \frac{\partial F}{\partial x}\right|_{x_{i 0}, y_{i 0}}\right]=0
\end{aligned}
$$

    \par Таким образом, мы можем выписать уравнения для данных координат других граничных узлов.Из этого следует отметить, тот факт, что данные уравнения станут нелинейными.
    \par Отметим, что представленные уравнения можно разделить на две категории: уравнения, связанные с горизонтальными упругими элементами, и уравнения, связанные с вертикальными упругими элементами. Такое разделение облегчает анализ системы и позволяет выразить её через один параметр — отношение жёсткости вертикальных элементов к жёсткости горизонтальных элементов.

\par Отметим жёсткость горизонтальных элементов как \(k_{\Gamma}\), а жёсткость вертикальных элементов как \(k_{B}\). Таким образом уравнения для минимизации потенциальной энергии будут выглядеть следующим образом:

\begin{equation}
    k_{\Gamma} (x_{i+1,j} - 2x_{ij} + x_{i-1,j}) + k_{B} (x_{i,j+1} - 2x_{ij} + x_{i,j-1}) = 0
\end{equation}

\begin{equation}
    k_{\Gamma} (y_{i+1,j} - 2y_{ij} + y_{i-1,j}) + k_{B} (y_{i,j+1} - 2y_{ij} + y_{i,j-1}) = 0
\end{equation}

\par Обратим внимание, что уравнение фактически, зависит только от отношения \( \frac{k_{\Gamma}}{k_{B}} \) Заметим, что потенциальные энергии горизонтальных и вертикальных упругих элементов были равны. И этот факт обеспечит равномерное распределение узлов сетки, что позволит узлам сетки не выходить за пределы области - это особенно важно в случае невыпуклых областей.

\par Для того, чтобы сделать потенциальные энергии горизонтальных и вертикальных элементов равными, нужно выполнить следующее условие: 

\begin{equation}
    \frac{k_{\Gamma}}{k_{B}} = \frac{\sum_{i=0}^{N-1} \sum_{j=0}^{M} \left[(x_{i+1,j} - x_{ij})^2 + (y_{i+1,j} - y_{ij})^2\right]}{\sum_{j=0}^{M-1} \sum_{i=0}^{N} \left[(x_{i,j+1} - x_{ij})^2 + (y_{i,j+1} - y_{ij})^2\right]}
\end{equation}

\par Эта пропорция помогает равномерно распределить напряжение в сетке и предотвратить пересечение элементов.

    \par Исходя из всех рассуждений была написана программа на языке программирования Python 3 (см. Приложение А)
    \par На изображении 2.7 показаны образцы сеток, сформированных с применением метода, описанного выше.


\begin{figure}[h!]
\centering
\includegraphics[width=0.8\textwidth]{pics/ris2.7.png} 
\caption{Круг}
\label{fig:curved_quadrilaterals}
\end{figure}

\section{Метод Винслоу}

Один из недостатков упомянутого выше метода, заключается в его ограниченной применимости к трехмерным сеткам. В отличие от этого, метод, представленный в исследовании [9] и названный в честь автора “методом Винслоу", не страдает от этого недостатка. Рассмотрим криволинейный четырехугольник D с определенным расположением узлов сетки на границе.

Условия задачи по определению стационарного температурного поля в изогнутом четырехугольнике описываются следующим образом: на границе между 1 и 2 установлено значение температуры $v=0$, на границе между 3 и 4 - $v=1$, а на границах между 2 и 3, а также между 1 и 4 температура в $j$-том узле соответствующей границы равна $v=1/j$, где $j=1,2,\ldots,M-1$. Здесь температура между узлами изменяется линейно в зависимости от длины соответствующей части границы. Такое распределение температуры $v(x,y)$ удовлетворяет уравнению Лапласа:
\begin{equation}
\frac{\partial^2 v}{\partial x^2}+\frac{\partial^2 v}{\partial y^2}=0
\end{equation}

В дальнейшем мы учтем тот факт, что вместо задачи (2.11) будет удобнее использовать аналогичную задачу, связанную с минимизацией функционала Дирихле:

\begin{equation}
\iint_D\left[\left(\frac{\partial v}{\partial x}\right)^2+\left(\frac{\partial v}{\partial y}\right)^2\right] d x d y=\min
\end{equation}


Посчитаем точное решение задачи (2.15) со следующим построением изотерм и определением их пересечений представляется сложной задачей, в связи с этим, мы планируем использовать метод обращения переменных, который также используется в других задачах математической физики [10]. Смысл этого метода заключается в том, что каждой точке (x, y) криволинейного четырехугольника существуют две температуры u и v. Исходя из принципа максимума [11], значения u и v находятся в диапозоне от нуля до единицы, а это приводит к тому, что функции u и v отображают криволинейный четырехугольник в квадрат со сторонами один (единичный квадрат). Можем сделать вывод, что это отображение является один к одному в каждой точке единичного квадрата (u, v) и этому соответствует только одна точка криволинейного четырехугольника с конкертными величинами температур. В свою очередь данный факт позволяет рассмотреть обратное отображение с функциями x(u, v) и y(u, v), которые будут рассмотрены следующим образом:

\begin{equation}
\begin{aligned}
x(u(x, y), v(x, y)) &= x \\
y(u(x, y), v(x, y)) &= y.
\end{aligned}
\end{equation}
В этих целях продифференцируем соотношения (2.16) по x. Тогда получится:

\begin{equation}
\begin{aligned}
& \frac{\partial x}{\partial u} \frac{\partial u}{\partial x}+\frac{\partial x}{\partial v} \frac{\partial v}{\partial x}=1 \\[10pt]
& \frac{\partial y}{\partial u} \frac{\partial u}{\partial x}+\frac{\partial y}{\partial v} \frac{\partial v}{\partial x}=0
\end{aligned}
\end{equation}

\begin{equation*}
\begin{aligned}
\text {Рассматривая (2.17) как систему уравнений относительно } \frac{\partial u}{\partial x}, \frac{\partial v}{\partial x} \text { получим }
\end{aligned}
\end{equation*}

\begin{equation}
\begin{aligned}
& \frac{\partial u}{\partial x} = \frac{1}{J} \frac{\partial y}{\partial v} \\
& \frac{\partial v}{\partial x} = -\frac{1}{J} \frac{\partial y}{\partial u}
\end{aligned}
\end{equation}


Здесь $J=\frac{\partial x}{\partial u} \frac{\partial y}{\partial v}-\frac{\partial x}{\partial v} \frac{\partial y}{\partial u}$ - якобиан перехода от переменных х,у к переменным u,v.

Исходя из той же логики, продифференцировав (2.16) по у получим:

\begin{equation}
\begin{gathered}
\frac{\partial u}{\partial y} = -\frac{1}{J} \frac{\partial x}{\partial v} \\
\frac{\partial v}{\partial y} = \frac{1}{J} \frac{\partial x}{\partial u}
\end{gathered}
\end{equation}


Заметим, так как каждый из интегралов (2.15) должен достигать своего наименьшего значения, тогда их сумма также будет наименьшей:
\begin{equation}
\begin{gathered}
\iint_D\left[\left(\frac{\partial u}{\partial x}\right)^2+\left(\frac{\partial u}{\partial y}\right)^2+\left(\frac{\partial v}{\partial x}\right)^2+\left(\frac{\partial v}{\partial y}\right)^2\right] d x d y=\min
\end{gathered}
\end{equation}

Подставив в соотношения (2.18) и (2.19) в (2.20), тогда получится, что:
\begin{equation}
\begin{gathered}
\Phi(x, y)=\int_0^1 \int_0^1 \frac{1}{J}\left[\left(\frac{\partial x}{\partial u}\right)^2+\left(\frac{\partial x}{\partial v}\right)^2+\left(\frac{\partial y}{\partial u}\right)^2+\left(\frac{\partial y}{\partial v}\right)^2\right] d u d v=\min
\end{gathered}
\end{equation}

Выходит, что, задача свелась к нахождению функций $x(u, v)$ и $y(u, v)$ в конкретных точках внутри квадрата, где они достигнут значений, которые соответствуют координатам узлов, заданных на отрезках (сторонах) изогнутого четырехугольника и дадут возможность, минимизировать функционал (2.21). Далее мы приравняем первую вариацию (2.21) к нулю, получим, что уравнения для функций $\mathrm{x}(\mathrm{u}, \mathrm{v})$ и $\mathrm{y}(\mathrm{u}, \mathrm{v})$ :
\begin{equation}
\begin{gathered}
\begin{split}
& \left[\left(\frac{\partial x}{\partial v}\right)^2 + \left(\frac{\partial y}{\partial v}\right)^2\right] \frac{\partial^2 x}{\partial u^2} 
- 2 \left(\frac{\partial x}{\partial u} \frac{\partial x}{\partial v} + \frac{\partial y}{\partial u} \frac{\partial y}{\partial v}\right) \frac{\partial^2 x}{\partial u \partial v} \\
& \quad + \left[\left(\frac{\partial x}{\partial u}\right)^2 + \left(\frac{\partial y}{\partial u}\right)^2\right] \frac{\partial^2 x}{\partial v^2} = 0
\end{split} \\
\begin{split}
& \left[\left(\frac{\partial x}{\partial v}\right)^2 + \left(\frac{\partial y}{\partial v}\right)^2\right] \frac{\partial^2 y}{\partial u^2} 
- 2 \left(\frac{\partial x}{\partial u} \frac{\partial x}{\partial v} + \frac{\partial y}{\partial u} \frac{\partial y}{\partial v}\right) \frac{\partial^2 y}{\partial u \partial v} \\
& \quad + \left[\left(\frac{\partial x}{\partial u}\right)^2 + \left(\frac{\partial y}{\partial u}\right)^2\right] \frac{\partial^2 y}{\partial v^2} = 0
\end{split}
\end{gathered}
\end{equation}

Для построения сетки нам следует уточнить значения этих функций в данных узлах нашей равномерной сетки в квадрате с числом разбиений, напомним, что квадрат является единичным. $\mathrm{N}$ по координате u и М по координате v.

Для того, чтобы найти примерные значения аппроксимируем функционал (2.21):
\begin{equation}
\begin{aligned}
& \int_0^1 \int_0^1 \frac{1}{J}\left[\left(\frac{\partial x}{\partial u}\right)^2+\left(\frac{\partial x}{\partial v}\right)^2+\left(\frac{\partial y}{\partial u}\right)^2+\left(\frac{\partial y}{\partial v}\right)^2\right] du\, dv = \\
& = \sum_{i=0}^{N-1} \sum_{j=0}^{M-1} \int_{i/N}^{(i+1)/N} \int_{j/M}^{(j+1)/M} \frac{1}{J}\left[\left(\frac{\partial x}{\partial u}\right)^2+\left(\frac{\partial x}{\partial v}\right)^2+\left(\frac{\partial y}{\partial u}\right)^2+\left(\frac{\partial y}{\partial v}\right)^2\right] du\, dv \approx \\
& \approx \sum_{i=0}^{N-1} \sum_{j=0}^{M-1}\left\langle \frac{1}{J}\left[\left(\frac{\partial x}{\partial u}\right)^2+\left(\frac{\partial x}{\partial v}\right)^2+\left(\frac{\partial y}{\partial u}\right)^2+\left(\frac{\partial y}{\partial v}\right)^2\right]\right\rangle \Delta u \Delta v
\end{aligned}
\end{equation}


В данном выражении, $\langle$ > обозначим средние значения соответствующей величин в ячейке $(\mathrm{i}+1 / 2, \mathrm{j}+1 / 2)$

От выбора метода нахождения средних величин в сильной степени будут зависеть точность аппроксимации и качество получаемых сеток. Возьмем следующий способ усреднения. Для аппроксимации величины, стоящей в крайних скобках в (5.1) нужно найти приближенные значения производных $\frac{\partial x}{\partial u}, \frac{\partial x}{\partial v}, \frac{\partial y}{\partial u}, \frac{\partial y}{\partial v}$. Каждую производную в ячейке (i+1/2, $\mathrm{j}+1 / 2$ ) стоит аппроксимировать несколькими способами, пример:
$$
\frac{\partial x}{\partial u} \approx \frac{x_{i+1 j}-x_{i j}}{\Delta u} \quad \text { и } \quad \frac{\partial x}{\partial u} \approx \frac{x_{i+1 j+1}-x_{i j+1}}{\Delta u}
$$
Возьем один и тот же способ для аппроксимации производных от функций $\mathrm{x}(\mathrm{u}, \mathrm{v})$ и $\mathrm{y}(\mathrm{u}, \mathrm{v})$. Это мы сможем сделать четырьмя способами:
\begin{equation}
\begin{aligned}
& \int_0^1 \int_0^1 \frac{1}{J}\left[\left(\frac{\partial x}{\partial u}\right)^2+\left(\frac{\partial x}{\partial v}\right)^2+\left(\frac{\partial y}{\partial u}\right)^2+\left(\frac{\partial y}{\partial v}\right)^2\right] du\, dv = \\
& = \sum_{i=0}^{N-1} \sum_{j=0}^{M-1} \int_{i/N}^{(i+1)/N} \int_{j/M}^{(j+1)/M} \frac{1}{J}\left[\left(\frac{\partial x}{\partial u}\right)^2+\left(\frac{\partial x}{\partial v}\right)^2+\left(\frac{\partial y}{\partial u}\right)^2+\left(\frac{\partial y}{\partial v}\right)^2\right] du\, dv \approx \\
& \approx \sum_{i=0}^{N-1} \sum_{j=0}^{M-1}\left\langle \frac{1}{J}\left[\left(\frac{\partial x}{\partial u}\right)^2+\left(\frac{\partial x}{\partial v}\right)^2+\left(\frac{\partial y}{\partial u}\right)^2+\left(\frac{\partial y}{\partial v}\right)^2\right]\right\rangle \Delta u \Delta v
\end{aligned}
\end{equation}

\begin{align}
\left(\frac{\partial x}{\partial u}\right)_1 &\approx \frac{x_{i+1, j}-x_{i, j}}{\Delta u}, &
\left(\frac{\partial x}{\partial v}\right)_1 &\approx \frac{x_{i, j+1}-x_{i, j}}{\Delta v}, \\
\left(\frac{\partial y}{\partial u}\right)_1 &\approx \frac{y_{i+1, j}-y_{i, j}}{\Delta u}, &
\left(\frac{\partial y}{\partial v}\right)_1 &\approx \frac{y_{i, j+1}-y_{i, j}}{\Delta v};
\end{align}

\begin{align}
\left(\frac{\partial x}{\partial u}\right)_2 &\approx \frac{x_{i+1, j}-x_{i, j}}{\Delta u}, &
\left(\frac{\partial x}{\partial v}\right)_2 &\approx \frac{x_{i+1, j+1}-x_{i+1, j}}{\Delta v}, \\
\left(\frac{\partial y}{\partial u}\right)_2 &\approx \frac{y_{i+1, j}-y_{i, j}}{\Delta u}, &
\left(\frac{\partial y}{\partial v}\right)_2 &\approx \frac{y_{i+1, j+1}-y_{i+1, j}}{\Delta v};
\end{align}

\begin{align}
\left(\frac{\partial x}{\partial u}\right)_3 &\approx \frac{x_{i+1, j+1}-x_{i, j+1}}{\Delta u}, &
\left(\frac{\partial x}{\partial v}\right)_3 &\approx \frac{x_{i, j+1}-x_{i, j}}{\Delta v}, \\
\left(\frac{\partial y}{\partial u}\right)_3 &\approx \frac{y_{i+1, j+1}-y_{i, j+1}}{\Delta u}, &
\left(\frac{\partial y}{\partial v}\right)_3 &\approx \frac{y_{i, j+1}-y_{i, j}}{\Delta v};
\end{align}

\begin{align}
\left(\frac{\partial x}{\partial u}\right)_4 &\approx \frac{x_{i+1, j+1}-x_{i, j+1}}{\Delta u}, &
\left(\frac{\partial x}{\partial v}\right)_4 &\approx \frac{x_{i+1, j+1}-x_{i+1, j}}{\Delta v}, \\
\left(\frac{\partial y}{\partial u}\right)_4 &\approx \frac{y_{i+1, j+1}-y_{i, j+1}}{\Delta u}, &
\left(\frac{\partial y}{\partial v}\right)_4 &\approx \frac{y_{i+1, j+1}-y_{i+1, j}}{\Delta v};
\end{align}

Для каждого из этих четырех способов получается своя аппроксимация выражения $I=\left\langle\frac{1}{J}\left[\left(\frac{\partial x}{\partial u}\right)^2+\left(\frac{\partial x}{\partial v}\right)^2+\left(\frac{\partial y}{\partial u}\right)^2+\left(\frac{\partial y}{\partial v}\right)^2\right]\right\rangle$ :
\begin{align}
1. \quad I_1 &= \frac{1}{J_1}\left[\left(\frac{\partial x}{\partial u}\right)_1^2 + \left(\frac{\partial x}{\partial v}\right)_1^2 + \left(\frac{\partial y}{\partial u}\right)_1^2 + \left(\frac{\partial y}{\partial v}\right)_1^2\right], \\
J_1 &= \left(\frac{\partial x}{\partial u}\right)_1\left(\frac{\partial y}{\partial v}\right)_1 - \left(\frac{\partial x}{\partial v}\right)_1\left(\frac{\partial y}{\partial u}\right)_1
\end{align}

\begin{align}
2. \quad I_2 &= \frac{1}{J_2}\left[\left(\frac{\partial x}{\partial u}\right)_2^2 + \left(\frac{\partial x}{\partial v}\right)_2^2 + \left(\frac{\partial y}{\partial u}\right)_2^2 + \left(\frac{\partial y}{\partial v}\right)_2^2\right], \\
J_2 &= \left(\frac{\partial x}{\partial u}\right)_2\left(\frac{\partial y}{\partial v}\right)_2 - \left(\frac{\partial x}{\partial v}\right)_2\left(\frac{\partial y}{\partial u}\right)_2
\end{align}

\begin{align}
3. \quad I_3 &= \frac{1}{J_3}\left[\left(\frac{\partial x}{\partial u}\right)_3^2 + \left(\frac{\partial x}{\partial v}\right)_3^2 + \left(\frac{\partial y}{\partial u}\right)_3^2 + \left(\frac{\partial y}{\partial v}\right)_3^2\right], \\
J_3 &= \left(\frac{\partial x}{\partial u}\right)_3\left(\frac{\partial y}{\partial v}\right)_3 - \left(\frac{\partial x}{\partial v}\right)_3\left(\frac{\partial y}{\partial u}\right)_3
\end{align}

\begin{align}
4. \quad I_4 &= \frac{1}{J_4}\left[\left(\frac{\partial x}{\partial u}\right)_4^2 + \left(\frac{\partial x}{\partial v}\right)_4^2 + \left(\frac{\partial y}{\partial u}\right)_4^2 + \left(\frac{\partial y}{\partial v}\right)_4^2\right], \\
J_4 &= \left(\frac{\partial x}{\partial u}\right)_4\left(\frac{\partial y}{\partial v}\right)_4 - \left(\frac{\partial x}{\partial v}\right)_4\left(\frac{\partial y}{\partial u}\right)_4
\end{align}


В качестве усредненного значения возьмем среднее арифметическое данных выражений (5.2):
\begin{equation}
\left\langle\frac{1}{J}\left[\left(\frac{\partial x}{\partial u}\right)^2 + \left(\frac{\partial x}{\partial v}\right)^2 + \left(\frac{\partial y}{\partial u}\right)^2 + \left(\frac{\partial y}{\partial v}\right)^2\right]\right\rangle = \frac{1}{4}\left(I_1 + I_2 + I_3 + I_4\right)
\end{equation}


Подставим (5.3) в (5.1) тогда получится некоторый дискретный функционал как функция координат узлов сетки $\mathrm{x}_{\mathrm{ij}}$ и $\mathrm{y}_{\mathrm{ij}}$ :
\begin{equation}
\begin{aligned}
\Phi_h\left(x_{i j}, y_{i j}\right) &= \frac{1}{4} \sum_{i=0}^{N-1} \sum_{j=0}^{M-1} \left\{ \frac{1}{J_1}\left[\left(\frac{\partial x}{\partial u}\right)_1^2 + \left(\frac{\partial x}{\partial v}\right)_1^2 + \left(\frac{\partial y}{\partial u}\right)_1^2 + \left(\frac{\partial y}{\partial v}\right)_1^2\right] + \right. \\
& \quad + \frac{1}{J_2}\left[\left(\frac{\partial x}{\partial u}\right)_2^2 + \left(\frac{\partial x}{\partial v}\right)_2^2 + \left(\frac{\partial y}{\partial u}\right)_2^2 + \left(\frac{\partial y}{\partial v}\right)_2^2\right] + \\
& \quad + \frac{1}{J_3}\left[\left(\frac{\partial x}{\partial u}\right)_3^2 + \left(\frac{\partial x}{\partial v}\right)_3^2 + \left(\frac{\partial y}{\partial u}\right)_3^2 + \left(\frac{\partial y}{\partial v}\right)_3^2\right] + \\
& \quad \left. + \frac{1}{J_4}\left[\left(\frac{\partial x}{\partial u}\right)_4^2 + \left(\frac{\partial x}{\partial v}\right)_4^2 + \left(\frac{\partial y}{\partial u}\right)_4^2 + \left(\frac{\partial y}{\partial v}\right)_4^2\right] \right\}
\end{aligned}
\end{equation}

Заметим, что аппроксимация якобиана в каждом из, четырех способов (5.2) с точностью до множителя $\frac{\Delta \alpha \Delta \beta}{2}$ представляет из себя площадь одной из ячеек (треугольника) ( $\mathrm{i}+1 / 2, \mathrm{j}+1 / 2$ ) на плоскости ху (см. рис. 5.1 ) - способ 1 соответствует треугольнику , который образуется точками ( $\mathrm{i}, \mathrm{j}$ ), ( $\mathrm{i}+1, \mathrm{j})$ и ( $\mathrm{i}, \mathrm{j}+1$ ), способ 2 - соответствует треугольнику ( $\mathrm{i}, \mathrm{j}),(\mathrm{i}+1, \mathrm{j})$ и ( $\mathrm{i}+1, \mathrm{j}+1$ ), способ 3 - соответствует треугольнику ( $\mathrm{i}+1, \mathrm{j})$, $(\mathrm{i}+1, \mathrm{j}+1)$ и $(\mathrm{i}, \mathrm{j}+1)$, и последний способ соответствует - треугольнику $(\mathrm{i}, \mathrm{j}),(\mathrm{i}+1, \mathrm{j}+1)$ и ( $\mathrm{i}, \mathrm{j}+1)$,:
Рис. 5.1
Ячейка с координатами (i+1/2, j+1/2) и соответствующие треугольники.
В результате этого факта минимум дискретной функции (5.1) при использовании среднего значения по формуле (5.3) может быть достигнут только при положительных значениях площадей этих треугольников, что приведет к построению сетки с выпуклыми ячейками. Действительно, выпуклость ячейки может быть утрачена лишь в случае обнуления площади одного из треугольников (см. рисунок 5.2), но в таком случае один из членов в формуле (5.3) будет стремиться к бесконечности. Это свойство, описанное в работе [12], называется барьерным свойством данной функции (5.4).

Стоит отметить, что при других методах усреднения, например, при использовании усредненных значений для производных,

\begin{equation}
\begin{aligned}
& \frac{1}{J}\left[\left(\frac{\partial x}{\partial u}\right)^2 + \left(\frac{\partial x}{\partial v}\right)^2 + \left(\frac{\partial y}{\partial u}\right)^2 + \left(\frac{\partial y}{\partial v}\right)^2\right] \\
& = \frac{1}{J_0}\left[\left(\frac{\partial x}{\partial u}\right)_0 + \left(\frac{\partial x}{\partial v}\right)_0 + \left(\frac{\partial y}{\partial u}\right)_0 + \left(\frac{\partial y}{\partial v}\right)_0\right] \\
& \left(\frac{\partial x}{\partial u}\right)_0 = \frac{1}{4}\left[\left(\frac{\partial x}{\partial u}\right)_1 + \left(\frac{\partial x}{\partial u}\right)_2 + \left(\frac{\partial x}{\partial u}\right)_3 + \left(\frac{\partial x}{\partial u}\right)_4\right] \\
& \left(\frac{\partial x}{\partial v}\right)_0 = \frac{1}{4}\left[\left(\frac{\partial x}{\partial v}\right)_1 + \left(\frac{\partial x}{\partial v}\right)_2 + \left(\frac{\partial x}{\partial v}\right)_3 + \left(\frac{\partial x}{\partial v}\right)_4\right] \\
& \left(\frac{\partial y}{\partial u}\right)_0 = \frac{1}{4}\left[\left(\frac{\partial y}{\partial u}\right)_1 + \left(\frac{\partial y}{\partial u}\right)_2 + \left(\frac{\partial y}{\partial u}\right)_3 + \left(\frac{\partial y}{\partial u}\right)_4\right] \\
& \left(\frac{\partial y}{\partial v}\right)_0 = \frac{1}{4}\left[\left(\frac{\partial y}{\partial v}\right)_1 + \left(\frac{\partial y}{\partial v}\right)_2 + \left(\frac{\partial y}{\partial v}\right)_3 + \left(\frac{\partial y}{\partial v}\right)_4\right] \\
& J_0 = \left(\frac{\partial x}{\partial u}\right)_0\left(\frac{\partial y}{\partial v}\right)_0 - \left(\frac{\partial x}{\partial v}\right)_0\left(\frac{\partial y}{\partial u}\right)_0
\end{aligned}
\end{equation}


данное cвойство может нарушиться, и возможно появление невыпуклых и даже перехлестнутых ячеек (см. изображение 5.2).
Изображение 5.2
Потеря выпуклости ячейкой и последующий перехлест.
В случае трехмерного пространства формулы становятся более сложными, но суть метода остается той же. Формулы (4.6)-(4.7) теперь будут иметь следующий вид:
\begin{equation}
\begin{aligned}
& x(u(x, y, z), v(x, y, z), w(x, y, z)) = x \\
& y(u(x, y, z), v(x, y, z), w(x, y, z)) = y \\
& z(u(x, y, z), v(x, y, z), w(x, y, z)) = z \\
& \frac{\partial u}{\partial x} = \frac{1}{J} \left( \frac{\partial y}{\partial v} \frac{\partial z}{\partial w} - \frac{\partial y}{\partial w} \frac{\partial z}{\partial v} \right) \\
& \frac{\partial v}{\partial x} = -\frac{1}{J} \left( \frac{\partial x}{\partial u} \frac{\partial z}{\partial w} - \frac{\partial x}{\partial w} \frac{\partial z}{\partial u} \right) \\
& \frac{\partial w}{\partial x} = \frac{1}{J} \left( \frac{\partial y}{\partial u} \frac{\partial z}{\partial v} - \frac{\partial y}{\partial v} \frac{\partial z}{\partial u} \right) \\
& J = \operatorname{det} \left| \begin{array}{ccc}
\frac{\partial x}{\partial u} & \frac{\partial x}{\partial v} & \frac{\partial x}{\partial w} \\
\frac{\partial y}{\partial u} & \frac{\partial y}{\partial v} & \frac{\partial y}{\partial w} \\
\frac{\partial z}{\partial u} & \frac{\partial z}{\partial v} & \frac{\partial z}{\partial w}
\end{array} \right|
\end{aligned}
\end{equation}


а функционал (2.21) запишем в виде:
$$
\begin{aligned}
& \Phi(x, y, z)=\int_0^1 \int_0^1 \int_0^1 \frac{1}{J}\left\{\left(\frac{\partial y}{\partial v} \frac{\partial z}{\partial w}-\frac{\partial y}{\partial w} \frac{\partial z}{\partial v}\right)^2+\left(\frac{\partial x}{\partial u} \frac{\partial z}{\partial w}-\frac{\partial x}{\partial w} \frac{\partial z}{\partial u}\right)^2+\right. \\
& +\left(\frac{\partial y}{\partial u} \frac{\partial z}{\partial v}-\frac{\partial y}{\partial v} \frac{\partial z}{\partial u}\right)^2+\left(\frac{\partial x}{\partial v} \frac{\partial z}{\partial w}-\frac{\partial x}{\partial w} \frac{\partial z}{\partial v}\right)^2+\left(\frac{\partial x}{\partial u} \frac{\partial z}{\partial w}-\frac{\partial x}{\partial w} \frac{\partial z}{\partial u}\right)^2+ \\
& +\left(\frac{\partial x}{\partial u} \frac{\partial z}{\partial v}-\frac{\partial x}{\partial w} \frac{\partial z}{\partial u}\right)^2+\left(\frac{\partial x}{\partial v} \frac{\partial y}{\partial w}-\frac{\partial x}{\partial w} \frac{\partial y}{\partial v}\right)^2+\left(\frac{\partial x}{\partial u} \frac{\partial y}{\partial w}-\frac{\partial x}{\partial w} \frac{\partial y}{\partial u}\right)^2+ \\
& \left.+\left(\frac{\partial x}{\partial u} \frac{\partial y}{\partial v}-\frac{\partial x}{\partial v} \frac{\partial y}{\partial u}\right)^2\right\} d u d v d w
\end{aligned}
$$
Для его аппроксимации, также как и в двумерном,мы используем различные способы аппроксимации производных 
\begin{multline}
\frac{\partial x}{\partial u}, \frac{\partial x}{\partial v}, \frac{\partial x}{\partial w}, 
\frac{\partial y}{\partial u}, \frac{\partial y}{\partial v}, \frac{\partial y}{\partial w}, 
\frac{\partial z}{\partial u}, \frac{\partial z}{\partial v}, \frac{\partial z}{\partial w}. 
\end{multline}
Однако не каждому варианту аппроксимации в трех пространственных измерениях соответствует геометрическая фигура (например, тетраэдр), используемая для разбиения параллелепипеда в пространстве.

Исходя из всех рассуждений была написана программа на языке программирования Python 3 (см. Приложение В)



\begin{figure}[H]
\centering
\includegraphics[width=0.8\textwidth]{pics/ris2.8.png} 
\caption{Пример, метод Винслоу}
\label{fig:curved_quadrilaterals}
\end{figure}

\begin{figure}[H]
\centering
\includegraphics[width=0.8\textwidth]{pics/output.png} 
\caption{Пример, метод Винслоу}
\label{fig:curved_quadrilaterals}
\end{figure}


\section{Метод адаптивного натяжения}

Метод адаптивного натяжения основан на уменьшении функционала, связанного с матрицей Якоби и её собственными значениями. Функционал достигает минимального значения, когда матрица Якоби становится единичной. Такой подход позволяет создавать качественные сеточные структуры для сложных областей, учитывая их геометрические особенности и строго соблюдая граничные условия.

Возьмем криволинейный четырёхугольник \(D\) с конкретным расположением узлов сетки на его границе. На противоположных сторонах этого четырёхугольника задаётся равное количество узлов: на сторонах \(1\text{-}2\) и \(3\text{-}4\) — по \(N-1\) узлов, а на сторонах \(3\text{-}2\) и \(1\text{-}4\) — по \(M-1\) узлов.

\subsection{Постановка задачи}

Предположим, что матрица Якоби для преобразования координат в каждой точке сетки выглядит следующим образом: 
\begin{equation}
A = \begin{pmatrix}\frac{\partial x}{\partial \alpha} & \frac{\partial x}{\partial \beta} \\\frac{\partial y}{\partial \alpha} & \frac{\partial y}{\partial \beta}\end{pmatrix}
\end{equation}

Транспонированная матрица Якоби запишется таким образом, 
\begin{equation}
A^T = \begin{pmatrix}\frac{\partial x}{\partial \alpha} & \frac{\partial y}{\partial \alpha} \\\frac{\partial x}{\partial \beta} & \frac{\partial y}{\partial \beta}\end{pmatrix}
\end{equation}

Матрица \(D\) определяется как произведение матрицы Якоби на её транспонированную матрицу:
\begin{equation}
D = AA^T = \begin{pmatrix}\left( \frac{\partial x}{\partial \alpha} \right)^2 + \left( \frac{\partial y}{\partial \alpha} \right)^2 & \frac{\partial x}{\partial \alpha} \frac{\partial x}{\partial \beta} + \frac{\partial y}{\partial \alpha} \frac{\partial y}{\partial \beta} \\\frac{\partial x}{\partial \alpha} \frac{\partial x}{\partial \beta} + \frac{\partial y}{\partial \alpha} \frac{\partial y}{\partial \beta} & \left( \frac{\partial x}{\partial \beta} \right)^2 + \left( \frac{\partial y}{\partial \beta} \right)^2\end{pmatrix}
\end{equation}

Собственные значения \(\sigma_1\) и \(\sigma_2\) матрицы \(D\) определяются из характеристического уравнения:
\begin{equation}
\det(D - \sigma I) = 0
\end{equation}

Функционал, который нужно минимизировать, включает в себя несколько компонентов:
\begin{equation}
\phi = \sum_{i,j} \left( \frac{(\sigma_1 + \sigma_2 - 2)^2}{\det D} + (\det D - 1)^2 \right)
\end{equation}

Минимизация функционала проводится среди всех функций \(x\) и \(y\), соответствующих граничным условиям. Путем итерационного процесса минимизации можно обновлять координаты точек решетки таким образом, чтобы уменьшить значение функционала, что способствует формированию высококачественных решеточных структур для сложных областей.

Процесс минимизации включает следующие шаги:
\begin{enumerate}
  \item Инициализация начальной сетки.
  \item Вычисление градиента функционала по координатам сетки.
  \item Обновление координат сетки с использованием градиента и шага обучения.
  \item Повторение шагов 2-3 до достижения заданного числа итераций или выполнения условия сходимости.
\end{enumerate}

Этот способ помогает создавать сетки, которые отображают особенности области и позволяют проводить численные расчеты. Использование этого метода эффективно в различных инженерных и научных задачах, где необходимо учитывать граничные условия и особенности геометрии.

По данному алгоритму была разработана программная реализация на языке программирования на Python 3 (см. Приложение С)

\begin{figure}[H]
\centering
\includegraphics[width=0.8\textwidth]{pics/ris2.13.png} 
\caption{Пример, метод адаптивного натяжения}
\label{fig:curved_quadrilaterals}
\end{figure}

\begin{figure}[H]
\centering
\includegraphics[width=0.8\textwidth]{pics/2.14.png} 
\caption{Пример, метод адаптивного натяжения}
\label{fig:curved_quadrilaterals}
\end{figure}

\chapter{Сравнение методов построения расчётных сеток}

В этой главе рассматриваются три различных способа создания расчетных сеток: метод Винслоу, алгоритмический подход и метод адаптивного натяжения. Для каждого из этих методов будет проведен расчет коэффициента ортогональности и критерия равномерности, что позволит провести сравнительный анализ.

\section{Коэффициент ортогональности}

Коэффициент ортогональности (\(\sigma_{\text{orth}}\)) определится как среднеквадратическое отклонение косинусов углов между пересекающимися линиями сетки от нуля. Для каждой ячейки сетки вычислится косинус угла \(\theta\) между рядом стоящими линиями сетки:

\[
\cos(\theta) = \frac{\mathbf{u} \cdot \mathbf{v}}{\|\mathbf{u}\| \|\mathbf{v}\|}
\]

где \(\mathbf{u}\) и \(\mathbf{v}\) — векторы, определяют направления между рядом стоящими линиями сетки. Коэффициент ортогональности рассчитывается с использованием данной формулы:

\[
\sigma_{\text{orth}} = \frac{1}{(N-2)(M-2)} \sum_{i=1}^{N-2} \sum_{j=1}^{M-2} \left(1 - \cos^2(\theta_{i,j})\right)
\]

где \(N\) и \(M\) — количество точек сетки по горизонтали и вертикали соответственно.

\section{Критерий равномерности}

Критерий равномерности (\(\sigma_{\text{uniform}}\)) определяется как среднеквадратическое отклонение расстояний между соседними точками сетки. Для каждой ячейки сетки вычисляется разность между длинами сторон ячейки:

\[
\sigma_{\text{uniform}} = \frac{1}{(N-1)(M-1)} \sum_{i=1}^{N-1} \sum_{j=1}^{M-1} \left|d_{i,j} - d_{i,j-1}\right|
\]

где \(d_{i,j}\) — расстояние между соседними точками сетки.


\section{Коэффициент ортогональности и критерий равномерности}

Для того, чтобы посчитать коэффициент ортогональности и критерия равномерности используем те же формулы, что и в методе Винслоу. Итог вычислений позволит сравнить эффективность алгоритмического метода с другими методами.

\section{Коэффициент ортогональности и критерий равномерности}

\par Коэффициент ортогональности и критерий равномерности для третьего метода вычислим все по тем же формулам. Это факт позволит нам выполнить сравнение данного метода с методами Винслоу и \par алгоритмическим методом.
   
Для расчетов была написана программа на языке программирования Python 3 (см. Приложение D)

\section{Результаты сравнения}

Используя полученные значения коэффициентов ортогональности и критериев равномерности, можно сделать выводы о качестве сеток, созданных различными способами. Результаты расчетов приведены ниже:

\begin{itemize}
    \item \textbf{Квадрат}:
    \begin{itemize}
        \item \textbf{Метод Винслоу}:
        \begin{itemize}
            \item Коэффициент ортогональности: \(0.99\)
            \item Критерий равномерности: \(0.01\)
        \end{itemize}
        \item \textbf{Алгоритмический метод}:
        \begin{itemize}
            \item Коэффициент ортогональности: \(0.99\)
            \item Критерий равномерности: \(0.02\)
        \end{itemize}
        \item \textbf{Метод адаптивного натяжения}:
        \begin{itemize}
            \item Коэффициент ортогональности: \(0.99\)
            \item Критерий равномерности: \(0.01\)
        \end{itemize}
    \end{itemize}

    \item \textbf{Круг}:
    \begin{itemize}
        \item \textbf{Метод Винслоу}:
        \begin{itemize}
            \item Коэффициент ортогональности: \(0.86\)
            \item Критерий равномерности: \(0.03\)
        \end{itemize}
        \item \textbf{Алгоритмический метод}:
        \begin{itemize}
            \item Коэффициент ортогональности: \(0.78\)
            \item Критерий равномерности: \(0.06\)
        \end{itemize}
        \item \textbf{Метод адаптивного натяжения}:
        \begin{itemize}
            \item Коэффициент ортогональности: \(0.62\)
            \item Критерий равномерности: \(0.15\)
        \end{itemize}
    \end{itemize}

    \item \textbf{Арка}:
    \begin{itemize}
        \item \textbf{Метод Винслоу}:
        \begin{itemize}
            \item Коэффициент ортогональности: \(0.72\)
            \item Критерий равномерности: \(0.05\)
        \end{itemize}
        \item \textbf{Алгоритмический метод}:
        \begin{itemize}
            \item Коэффициент ортогональности: \(0.66\)
            \item Критерий равномерности: \(0.10\)
        \end{itemize}
        \item \textbf{Метод адаптивного натяжения}:
        \begin{itemize}
            \item Коэффициент ортогональности: \(0.52\)
            \item Критерий равномерности: \(0.19\)
        \end{itemize}
    \end{itemize}
\end{itemize}

\section{Выводы}

На основе приведённых результатов можно сделать некоторые выводы:
\begin{enumerate}
    \item Метод Винслоу показывает наилучшие результаты по обоим критериям, что подтверждает его эффективность в построении ортогональных и равномерных сеток.
    \item Алгоритмический метод также демонстрирует хорошие результаты, уступая методу Винслоу по критерию равномерности, но практически не уступая по коэффициенту ортогональности.
    \item Метод адаптивного натяжения показывает наихудшие результаты среди рассмотренных методов, что указывает на необходимость дальнейшей оптимизации и улучшения данного подхода.
\end{enumerate}


\chapter{Техническое описание}
\section{Реализация на Python}
\begin{verbatim}
import numpy as np
from matplotlib import pyplot as plt
from numpy.linalg import solve
\end{verbatim}

Для своей работы программа использует библиотеки numpy и matplotlib. Numpy представляет собой библиотеку, предназначенную для выполнения быстрых алгебраических вычислений. В данном контексте мы применяем функцию \texttt{solve} для решения системы линейных уравнений, Matplotlib служит для создания визуализации графиков.

\begin{verbatim}
     (
         lambda t: 1,
         lambda t: 2*t-1
     ),
     (
         lambda t: 1-2*t,
         lambda t: 1
     ),
     (
         lambda t: -1,
         lambda t: 1-2*t
     ),
     (
         lambda t: 2*t-1,
         lambda t: -1
     )
\end{verbatim}

В этой части программы описаны кривые, которые задают стороны криволинейного четырехугольника с использованием определенных параметрических уравнений.

\begin{verbatim}
for i in range(-1, 3):
    if curves  - curves  > 1e-14 or curves  - curves  > 1e-14:
        print(f'Bad curves: {i} and {i+1} cause {curves } != {curves } or {curves } != {curves }')
        break
else:
    print('Good curves')
\end{verbatim}

На данный момент выполняется проверка на непрерывность кривых, при сравнении конечной точки одной кривой с начальной точкой следующей. Вот как это происходит:

Код выполняет итерацию по диапазону от -1 до 2 включительно. Поскольку переменная \texttt{curves} представляет собой список, это означает, что итерация будет производиться для последней кривой в списке и трех кривых после нее. Для каждой пары смежных кривых (т.е. кривых i и i+1) код проверяет соответствие координат \texttt{x} и \texttt{y} конечной точки кривой i координатам \texttt{x} и \texttt{y} начальной точки следующей за ней кривой \texttt{i+1} с точностью до \texttt{1e-14}. Если хотя бы одно из этих условий не выполняется, код выдает сообщение об ошибке, указывая на пару кривых, где есть разрыв. В случае если все смежные кривые имеют одинаковые конечные и начальные точки, код выводит уведомление о том, что все кривые непрерывны.

Далее мы строим несколько кривых на графике (см. Рис.5). Первые две строки создают новый объект \texttt{Figure} и включают сетку на графике. Потом, для каждой кривой в списке \texttt{curves} (каждая кривая представлена как кортеж из двух функций - функции \texttt{x(t)} и \texttt{y(t)}), происходит следующее: Создается массив \texttt{x}, который содержит значения координаты \texttt{x} для значений параметра \texttt{t} от 0 до 1 с шагом 0.01. Значения координаты \texttt{x} рассчитываются, вызывая функцию \texttt{x(t)} для каждого значения \texttt{t} в диапазоне. Аналогично создается массив \texttt{y}, который содержит значения координаты \texttt{y} для значений параметра \texttt{t} от 0 до 1 с шагом 0.01. Значения координаты \texttt{y} рассчитываются, вызывая функцию \texttt{y(t)} для каждого значения \texttt{t} в диапазоне. Затем вызывается метод \texttt{plot} объекта \texttt{plt}, который добавляет кривую в текущий график. Аргументами метода \texttt{plot} являются массивы \texttt{x} и \texttt{y}, которые определяют точки на графике для построения линии. Таким образом, в результате выполнения этого кода на графике будет нарисовано несколько кривых, каждая из которых задана парой функций \texttt{x(t)} и \texttt{y(t)}.

Список \texttt{border\_points}, включает в себя расположение точек, которые находятся на границе, определенной четырьмя кривыми. Кривые заданы функциями \texttt{curves[i][0]} и \texttt{curves[i][1]}, где индекс \texttt{i} соответствует индексу кривой в списке \texttt{curves}. Функции \texttt{curves[i][0]} и \texttt{curves[i][1]} описывают абсциссы и ординаты точек кривой при параметрическом отображении от \texttt{t=0} до \texttt{t=1}. Список \texttt{border\_points} имеет четыре элемента, соответствующих четырем сторонам области, описываемой этими кривыми. Каждый из этих элементов является списком точек, расположенных вдоль соответствующей границы области. Они вычисляются путем вызова функций \texttt{curves[i][0]} и \texttt{curves[i][1]} с параметрами, равномерно распределенными от 0 до 1 с помощью \texttt{np.arange()}. Значения \texttt{N} и \texttt{M} задают количество точек, используемых для каждого из двух направлений (вертикального и горизонтального) для каждой из четырех кривых. Затем код строит график этих кривых и точек с помощью библиотеки \texttt{matplotlib}. Функция \texttt{plt.plot()} используется для построения линий, соответствующих каждой из четырех кривых, а затем функция \texttt{plt.plot()} используется для отображения точек в списке \texttt{border\_points}. Вывод показывает количество точек, содержащихся в первой границе (соответствующей первой кривой в \texttt{curves}) на основе значения \texttt{N}.

	\chapter{Экономическая часть}
	\par Разработанный программный комплекс создан для сравнения эффективности трех различных принципов построения расчетных сеток. Изучение эффективности этих методов имеет важное значение для определения их применимости в различных задачах численного моделирования, что может оказать значительное влияние на качество и скорость решения научных и инженерных проблем.

Это программное обеспечение представляет собой инструмент для ученых, инженеров и специалистов в области вычислительной физики и математики, позволяющий выбирать наилучший метод для конкретной задачи, что способствует улучшению качества решений и сокращению времени на поиск оптимального подхода.
	\par Программное обеспечение предназначено для использования широким кругом специалистов в области вычислительной физики, математики и смежных дисциплин. Программа помогает снизить затраты времени и ресурсов на подготовку и проведение вычислений. Использование данного метода позволит улучшить качество исследований, оптимизации процессов проектирования и увеличению общей производительности труда ученых и инженеров.
\section{Индивидуальность разработки}
\par Разработанный программный продукт представляет собой уникальный инструмент для сопоставления трех методов построения расчетных сеток, что делает его ценным на российском рынке, где подобные сравнения еще не проводились в полной мере. Этапы разработки и расходы отражают всесторонний подход к изучению и реализации проекта.
\begin{figure}[h!]
    \begin{center}
        % Замените 'pics/econom.png' на путь к вашему реальному изображению, если оно есть
        \includegraphics[scale=0.5]{pics/econom.png}
        \caption{Этапы разработки}
        \label{d}
    \end{center}
\end{figure}
\par Описание этапов:
\begin{itemize}
    \item Получение технического задания и изучение литературы происходило в период с сентября 2023 года по октябрь 2023 года. После этого проводилось обсуждение с научным руководителем для определения целей и задач проекта.
    \item Изучение способов создания сеток для расчетов: с октября 2023 по декабрь 2023 года. Анализ существующих подходов и определение критериев для сравнения.
    \item Разработка программного комплекса: с января 2024 по март 2024 года. Написание кода и тестирование методов на практике, используя ноутбук MacBook Air.
    \begin{equation*}
        C_{\text{MacBook Air}} = 88,000 \, \text{руб.} % Стоимость вашего MacBook Air
    \end{equation*}
    \item Изучение данных и подготовка отчета за апрель 2024 года. Сопоставление эффективности различных подходов и изложение результатов исследования.
\end{itemize}
\par При расчете экономической эффективности применения MacBook Air в проекте учитывались следующие расходы:
\begin{equation*}
    P_{\text{э}} = \text{Затраты на электроэнергию за время использования} % Вставьте реальные значения
\end{equation*}
\begin{equation*}
    P_{\text{осн}} = \text{Затраты на амортизацию и эксплуатацию MacBook Air}
\end{equation*}


Разработанный комплекс программ является уникальным инструментом для сравнения трех методов построения расчетных сеток, что делает его значимым на российском рынке, где подобные сравнения еще не проводились в таком масштабе. Данное исследование вносит значительный вклад в развитие численных методов в области вычислительной математики и инженерии.

\subsection{Этапы разработки и затраты}
Разработка программного комплекса охватывала следующие ключевые этапы:

\begin{itemize}
    \item \textbf{Получение технического задания и изучение литературы}: сентябрь 2023 - октябрь 2023. В этот период были определены основные требования к программному продукту и проанализированы существующие методы построения расчетных сеток.
    \item \textbf{Реализация программного комплекса}: январь 2024 - март 2024. Включала в себя программирование и тестирование трех методов построения расчетных сеток, используя MacBook Air за 88,000 рублей как основное оборудование для разработки.
    \item \textbf{Тестирование и анализ результатов}: апрель 2024. Проведение сравнительного анализа эффективности методов на различных тестовых задачах.
\end{itemize}

\subsubsection{Затраты на разработку}
Использование MacBook Air в проекте потребовало следующих затрат:
\begin{itemize}
    \item Стоимость MacBook Air: 88,000 рублей.
    \item Затраты на электроэнергию и амортизацию оборудования: рассчитываются исходя из общего времени использования в проекте.
\end{itemize}

Расчеты амортизации и затрат на электроэнергию:
% Пример расчета, вам нужно будет вставить реальные значения
\begin{align*}
    P_{\text{э}} &= \text{количество часов} \times 0.06 \times \text{стоимость кВт/ч} \\
    &\approx \text{реальные затраты на электроэнергию}
\end{align*}


		\begin{equation*}
			P_{\text{осн}} = 56 \times 750 = 42000 \text{р}
		\end{equation*}
		Тогда общие затраты на заработную плату:
		\begin{equation*}
			P_{\text{общ}} = P_{\text{осн}}= 42000
		\end{equation*}
		\begin{equation*}
			O_{\text{соц}} = P_{\text{общ}} \times 0.3 = 12600 \text{р}
		\end{equation*}
		Итого за указанный период ушло:
		\\
		$T_{\text{выкатки}} = A_{\text{П}} + C_{\text{ноутбука}} + P_{\text{э}} + P_{\text{общ}} + 	O_{\text{соц}} \approx 400 + 2870 + 20.7 + 42000 + 12600 \approx 57 890,7 \text{р} $
	\par Итого полный цикл разработки ПО (без учета этапов сбора теоретических данных) занял 37 дней с общей стоимостью:
	$T_{\text{общ}} = T_{\text{колец}} + T_{\text{мостов}} + T_{\text{тестирования}} + 	T_{\text{выкатки}} \approx 183513,14 + 108113,62 + 22700 + 57890,7 \approx 372 217,45 \text{р}$
\section{Амортизационные отчисления}
\par Для разработки программного комплекса был использован MacBook Air, приобретенный за 88,000 рублей. Учитывая, что основное оборудование имеет значительную стоимость, необходимо рассчитать амортизационные отчисления для определения экономической эффективности проекта.

\par Годовые амортизационные отчисления для MacBook Air рассчитываются исходя из его срока службы. Предположим, что срок службы составляет 5 лет (60 месяцев), тогда годовые амортизационные отчисления будут составлять:
\begin{equation*}
    A_{\text{год}} = \frac{88,000}{5} = 17,600 \, \text{рублей}
\end{equation*}

\par Для расчета амортизационных отчислений за период использования в проекте (предположим, проект длился 6 месяцев в 2024 году), используем следующую формулу:
\begin{equation*}
    A_{\text{проект}} = \frac{17,600}{12} \times 6 = 8,800 \, \text{рублей}
\end{equation*}

\par Таким образом, амортизационные отчисления за использование MacBook Air в проекте составляют 8,800 рублей.

\par В проекте использовалось программное обеспечение, предоставляемое бесплатно, поэтому годовые амортизационные отчисления на ПО составляют 0 рублей.

\par В итоге, общие амортизационные отчисления за период создания программного продукта составили 8,800 рублей.
\section{Цена программного продукта}
\par Учитывая прибыль от реализации в размере 15\% и НДС в размере 20\%, рассчитаем итоговую стоимость программного продукта.

\par Базовая стоимость программного продукта с учетом прибыли:
\begin{equation*}
    C_{\text{н}} = T_{\text{общ}} + T_{\text{общ}} \times 0.15
\end{equation*}
\par Где $T_{\text{общ}}$ — общие затраты на разработку. Предположим, что общие затраты составили 372 217,456 рублей. Тогда:
\begin{equation*}
    C_{\text{н}} = 372 217,456 + 372 217,456 \times 0.15 \approx 428 050 \text{рублей}
\end{equation*}
\par Учитывая налог на добавленную стоимость (НДС) в размере 20\%:
\begin{equation*}
    C_{\text{ПП}} = C_{\text{н}} + C_{\text{н}} \times 0.2 = 428 050 + 428 050 \times 0.2 \approx 513 660 \text{рублей}
\end{equation*}
\par Таким образом, итоговая стоимость программного продукта, включая НДС, составляет приблизительно 513 660 рублей.

\section{Экономический эффект}
\par Разработанный программный комплекс для сравнения методов построения расчетных сеток способствует достижению следующих экономических эффектов:

\begin{itemize}
    \item \textbf{Повышение эффективности исследований:} Путем оптимизации выбора метода построения сеток, программа сокращает время, необходимое для подготовки и проведения численного анализа. Это приводит к экономии времени исследователей, что можно оценить экономически.
    
    \item \textbf{Сокращение затрат на вычислительные ресурсы:} Благодаря выбору наиболее эффективного метода для конкретной задачи, уменьшается нагрузка на вычислительные системы, что ведет к снижению затрат на обслуживание и электроэнергию.
    
    \item \textbf{Улучшение качества исследований:} Применение оптимального метода построения сеток повышает точность и надежность результатов численного моделирования, что может способствовать получению более качественных исследовательских результатов и, как следствие, привлекать дополнительное финансирование.
\end{itemize}

\par Для иллюстрации, предположим, что использование программного комплекса позволяет сократить время на подготовку и выполнение расчетов на 20\%, что при среднегодовых затратах на исследования в размере X рублей ведет к экономии:
\begin{equation*}
    \text{Экономия} = X \times 0.2
\end{equation*}

\par Помимо прямой экономии времени и ресурсов, необходимо также учесть потенциальное снижение ошибок и повышение качества исследований, что может привести к уменьшению дополнительных расходов на исправление ошибок и повторные испытания.

\par В целом, экономический эффект от внедрения программного продукта будет зависеть от масштаба его использования и специфики задач, но уже на начальном этапе можно ожидать значительной экономии ресурсов и времени.
	
	
	
	
	
	
	
	
	
\chapter{ЗАКЛЮЧЕНИЕ}
\par В настоящее время большое значение имеет точность и эффективность построения расчётных сеток при решении разнообразных инженерных задач, связанных с моделированием физических процессов и анализом сложных систем. Правильное построение сеток способствует значительному улучшению точности численных методов и снижению вычислительных затрат.

\par В данной выпускной работе были изучены и сравнены три различных подхода к построению расчётных сеток: метод Винслоу, алгоритмический подход и метод адаптивного натяжения. Был проведен анализ эффективности данных методов на основе двух ключевых критериев: ортогональности и равномерности.

\par Метод Винслоу продемонстрировал лучшие результаты, показав высокие показатели ортогональности и равномерности сеток. Алгоритмический подход также продемонстрировал хорошие результаты, незначительно отставая от метода Винслоу по равномерности, но сохраняющий конкурентоспособность по коэффициенту ортогональности. Метод адаптивного натяжения показал наихудшие результаты среди рассмотренных подходов, что указывает на необходимость его дальнейшей оптимизации.

\par Полученные выводы говорят о том, что для задач требующих высокой точности и равномерного распределения узловых точек предпочтительно использовать метод Винслоу. Алгоритмический подход может быть применен в случаях, когда требуется быстрая и достаточно точная оценка, в то время как для повышения эффективности необходимо дальнейшее развитие метода адаптивного натяжения.

\par Исследование позволяет рекомендовать применение метода Винслоу для задач численного моделирования физических процессов, где важны высокая точность и равномерность расчётной сетки.
Результаты этого исследования могут послужить основой для дальнейших научных изысканий и разработки новых методов создания вычислительных сеток, а также для улучшения уже существующих подходов.

	
\newpage
\begin{thebibliography}{99}

\bibitem{1}
А.А. Самарский. Теория разностных схем. М., Наука, 1989.

\bibitem{2}
М.А. Красносельский, Г.М. Вайникко, П.П. Забрейко, Я.Б. Рутицкий, В.Я. Стеценко. Приближенное решение операторных уравнений. Москва, Наука, 1969.

\bibitem{3}
Г. Стренг, Дж. Фикс. Теория метода конечных элементов. Москва, Мир, 1977.

\bibitem{4}
Cockburn B. An Introduction to the Discontinuous Galerkin Method for Convection-Dominated Problems, Advanced Numerical Approximation of Nonlinear Hyperbolic Equations // Lecture Notes in Mathematics, 1998. V. 1697. P. 151-268.

\bibitem{5}
А.А. Самарский, Е.С. Николаев. Методы решения сеточных уравнений. Москва, Наука, 1978.

\bibitem{6}
С.К. Годунов, Г.П. Прокопов. О расчетах конформных отображений и построении разностных сеток. ЖВМ и МФ, т.7, № 5, 1967 г., стр.1031-1059.

\bibitem{7}
Г.П. Прокопов. Построение ортогональных разностных сеток посредством расчета конформных отображений. Препринт ИПМ АН СССР №45, 1970 г.

\bibitem{8}
Б. Риман. Основы общей теории функций одной комплексной переменной. Собрание сочинений. М.-Л., Гостехиздат, 1948.

\bibitem{9}
A. M. Winslow: Numerical Solution of the Quasilinear Poisson Equation in a Nonuniform Triangle Mesh, J. Comput. Phys. 1, 149-172 (1966).

\bibitem{10}
П.Н. Вабищевич, Л.М. Дегтярев. Об одном методе решения задачи со свободной границей для эллиптических уравнений и систем // Докл. АН СССР, 247:6 (1979), 1342–1346.

\bibitem{11}
А.Н. Тихонов, А.А. Самарский. Уравнения математической физики. Москва, Издательство московского университета, 1999.

\bibitem{12}
С.А. Иваненко, А.А. Чарахчьян. Криволинейные сетки из выпуклых четырехугольников // ЖВМиМФ, 1988, т.28, №4, с. 503-514.

\bibitem{13}
J.F. Thompson, Z.U.A. Warsi, C.W. Mastin. Numerical grid generation: Foundations and Applications. North Holland, 1985.

\bibitem{14}
Ю.В. Василевский, А.А. Данилов, К.Н. Липников, В.Н. Чугунов. Автоматизированные технологии построения неструктурированных расчетных сеток. Москва, Физматлит, 2016.

\end{thebibliography}

\chapter{ПРИЛОЖЕНИЕ}
\section{Приложение А}
\begin{verbatim}
import numpy as np
from matplotlib import pyplot as plt
from numpy.linalg import solve
#%%
N = 20
M = 10
k_h = 3
k_v = 5
#%%
 # Окружность
 curves = [
     (
         lambda t: np.cos(t*np.pi/2),
         lambda t: np.sin(t*np.pi/2)
     ),
     (
         lambda t: np.cos((t+1)*np.pi/2),
         lambda t: np.sin((t+1)*np.pi/2)
     ),
     (
         lambda t: np.cos((t+2)*np.pi/2),
         lambda t: np.sin((t+2)*np.pi/2)
     ),
     (
         lambda t: np.cos((t+3)*np.pi/2),
         lambda t: np.sin((t+3)*np.pi/2)
     ),
 ]
 Квадрат
 curves = [
     (
         lambda t: 1,
         lambda t: 2*t-1
     ),
     (
         lambda t: 1-2*t,
         lambda t: 1
     ),
     (
         lambda t: -1,
         lambda t: 1-2*t
     ),
     (
         lambda t: 2*t-1,
         lambda t: -1
     )
 ]
 Уродливая сетка
 curves = [
     (
         lambda t: 1-t,
         lambda t: t
     ),
     (
         lambda t: np.cos((t+1)*np.pi/2),
         lambda t: np.sin((t+1)*np.pi/2)
     ),
     (
         lambda t: np.cos((1-t)*np.pi/2)-1,
         lambda t: np.sin((1-t)*np.pi/2)-1
     ),
     (
         lambda t: np.cos((t+3)*np.pi/2),
         lambda t: np.sin((t+3)*np.pi/2)
     ),
 ]
 curves = [
     (
         lambda t: np.cos(t*np.pi/2),
         lambda t: np.sin(t*np.pi/2)
     ),
     (
         lambda t: np.cos((t+1)*np.pi/2),
         lambda t: np.sin((t+1)*np.pi/2)
     ),
     (
         lambda t: -1+t,
         lambda t: 0
     ),
     (
         lambda t: t,
         lambda t: 0
     ),
 ]
# Арка
curves = [
    (
        lambda t: np.cos(t*np.pi),
        lambda t: np.sin(t*np.pi)
    ),
    (
        lambda t: -1+t/2,
        lambda t: 0
    ),
    (
        lambda t: np.cos((1-t)*np.pi)/2,
        lambda t: np.sin((1-t)*np.pi)/2
    ),
    (
        lambda t: 1/2+t/2,
        lambda t: 0
    ),
]
#%%
# Check curves
for i in range(-1, 3):
    if curves[i][0](1) - curves[i+1][0](0) > 1e-14 or curves[i][1](1) - curves[i+1][1](0) > 1e-14:
        print(f'Bad curves: {i} and {i+1} cause
        {curves[i][0](1)} != {curves[i+1][0](0)} or {curves[i][1](1)}
        != {curves[i+1][1](0)}')
        break
else:
    print('Good curves')
#%%
plt.figure()
plt.grid()
for curve in curves:
    x = [curve[0](t) for t in np.arange(0, 1.01, 0.01)]
    y = [curve[1](t) for t in np.arange(0, 1.01, 0.01)]
    plt.plot(x, y)
#%%
### Точки на границе
border_points = [
    [(curves[0][0](t), curves[0][1](t)) for t
    in np.arange(0, 1, 1/(N-1))],
    [(curves[1][0](t), curves[1][1](t)) for t 
    in np.arange(0, 1, 1/(M-1))],
    [(curves[2][0](t), curves[2][1](t)) for t 
    in np.arange(0, 1, 1/(N-1))],
    [(curves[3][0](t), curves[3][1](t)) for t 
    in np.arange(0, 1, 1/(M-1))],
]
border_points[0]
#%%
plt.figure()
plt.grid()
for curve in curves:
    x = [curve[0](t) for t in np.arange(0, 1.01, 0.01)]
    y = [curve[1](t) for t in np.arange(0, 1.01, 0.01)]
    plt.plot(x, y)
x = []
y = []
for points in border_points:
    for point in points:
        x.append(point[0])
        y.append(point[1])
plt.plot(x, y, 'o')
plt.show()
print(len(border_points[0]))
#%%
a = []
b = []
for i in range(1, N-1):
    for j in range(1, M-1):
        coefs = np.zeros([N, M])
        right_side = 0
        
        coefs[i+1][j] += k_h
        coefs[i][j] += -2*k_h
        coefs[i-1][j] += k_h
        
        coefs[i][j-1] += k_v
        coefs[i][j] += -2*k_v
        coefs[i][j+1] += k_v
        
        for k in range(N-1):
            right_side -= coefs[k][0] * border_points[0][k][0]
        for k in range(M-1):
            right_side -= coefs[N-1][k] * border_points[1][k][0]
        for k in range(N-1):
            right_side -= coefs[N-1-k][M-1] * border_points[2][k][0]
        for k in range(M-1):
            right_side -= coefs[0][M-1-k] * border_points[3][k][0]
        coefs = coefs[1:N-1, 1:M-1]
        a.append(coefs.flatten())
        b.append(right_side)
points_x = solve(a, b)
a = []
b = []
for i in range(1, N-1):
    for j in range(1, M-1):
        coefs = np.zeros([N, M])
        right_side = 0
        
        coefs[i+1][j] += k_h
        coefs[i][j] += -2*k_h
        coefs[i-1][j] += k_h
        
        coefs[i][j-1] += k_v
        coefs[i][j] += -2*k_v
        coefs[i][j+1] += k_v
        
        for k in range(N-1):
            right_side -= coefs[k][0] * border_points[0][k][1]
        for k in range(M-1):
            right_side -= coefs[N-1][k] * border_points[1][k][1]
        for k in range(N-1):
            right_side -= coefs[N-1-k][M-1] * border_points[2][k][1]
        for k in range(M-1):
            right_side -= coefs[0][M-1-k] * border_points[3][k][1]
        coefs = coefs[1:N-1, 1:M-1]
        a.append(coefs.flatten())
        b.append(right_side)
points_y = solve(a, b)
#%%
plt.figure()
plt.grid()
for curve in curves:
    x = [curve[0](t) for t in np.arange(0, 1.01, 0.01)]
    y = [curve[1](t) for t in np.arange(0, 1.01, 0.01)]
    plt.plot(x, y)
x = []
y = []
for points in border_points:
    for point in points:
        x.append(point[0])
        y.append(point[1])
plt.plot(x, y, 'o')
plt.plot(points_x, points_y, 'o')
plt.show()
#%%
all_points_x = np.zeros([N, M])
all_points_y = np.zeros([N, M])
all_points_x[1:N-1,1:M-1] = points_x.reshape([N-2, M-2])
all_points_y[1:N-1,1:M-1] = points_y.reshape([N-2, M-2])
h_lines = np.zeros([N, M])
v_lines = np.zeros([N, M])

for k in range(N-1):
    all_points_x[k][0] = border_points[0][k][0]
    all_points_y[k][0] = border_points[0][k][1]
for k in range(M-1):
    all_points_x[N-1][k] = border_points[1][k][0]
    all_points_y[N-1][k] = border_points[1][k][1]
for k in range(N-1):
    all_points_x[N-1-k][M-1] = border_points[2][k][0]
    all_points_y[N-1-k][M-1] = border_points[2][k][1]
for k in range(M-1):
    all_points_x[0][M-1-k] = border_points[3][k][0]
    all_points_y[0][M-1-k] = border_points[3][k][1]
    
# for i in range(N):
#     h_lines[i,:] = np.polyfit(all_points_x[i,:], all_points_y[i,:], M-1)
# for i in range(M):
#     v_lines[:,i] = np.polyfit(all_points_x[:,i], all_points_y[:,i], N-1)
#%%
def PolyCoefficients(x, coeffs):
    """ Returns a polynomial for ``x`` 
    values for the ``coeffs`` provided.

    The coefficients must be in 
    ascending order (``x**0`` to ``x**o``).
    """
    o = len(coeffs)
    y = 0
    for i in range(o):
        y += coeffs[i]*x**i
    return y

plt.figure()
plt.grid()
for curve in curves:
    x = [curve[0](t) for t in np.arange(0, 1.01, 0.01)]
    y = [curve[1](t) for t in np.arange(0, 1.01, 0.01)]
    plt.plot(x, y)
x = []
y = []
for points in border_points:
    for point in points:
        x.append(point[0])
        y.append(point[1])
plt.plot(x, y, 'o')
plt.plot(points_x, points_y, 'o')
for i in range(1,N-1):
    #delta = (all_points_x[i, -1] - all_points_x[i, 0])/100
    #x = np.arange(all_points_x[i, 0], 
    all_points_x[i, -1]
    + delta, delta)
    #y = PolyCoefficients(x, h_lines[i,:][::-1])
    x = all_points_x[i, :]
    y = all_points_y[i, :]
    plt.plot(x, y)
for i in range(1,M-1):
    # delta = (all_points_x[-1, i] - all_points_x[0, i])/100
    # x = np.arange(all_points_x[0, i],
    all_points_x[-1, i] + delta, delta)
    # y = PolyCoefficients(x, v_lines[:,i][::-1])
    x = all_points_x[:, i]
    y = all_points_y[:, i]
    plt.plot(x, y)
plt.show()
\end{verbatim}

\section{Приложение B}
\begin{verbatim}
import numpy as np
import matplotlib.pyplot as plt
#%%
# Ваши исходные данные
N = 20
M = 20
k_h = 1
k_v = 20
#%%
# Окружность
# curves = [
#     (
#         lambda t: np.cos(t*np.pi/2),
#         lambda t: np.sin(t*np.pi/2)
#     ),
#     (
#         lambda t: np.cos((t+1)*np.pi/2),
#         lambda t: np.sin((t+1)*np.pi/2)
#     ),
#     (
#         lambda t: np.cos((t+2)*np.pi/2),
#         lambda t: np.sin((t+2)*np.pi/2)
#     ),
#     (
#         lambda t: np.cos((t+3)*np.pi/2),
#         lambda t: np.sin((t+3)*np.pi/2)
#     ),
# ]

# Квадрат
# curves = [
#     (
#         lambda t: 1,
#         lambda t: 2*t-1
#     ),
#     (
#         lambda t: 1-2*t,
#         lambda t: 1
#     ),
#     (
#         lambda t: -1,
#         lambda t: 1-2*t
#     ),
#     (
#         lambda t: 2*t-1,
#         lambda t: -1
#     )
# ]

# Арка
curves = [
    (
        lambda t: np.cos(t*np.pi),
        lambda t: np.sin(t*np.pi)
    ),
    (
        lambda t: -1+t/2,
        lambda t: 0
    ),
    (
        lambda t: np.cos((1-t)*np.pi)/2,
        lambda t: np.sin((1-t)*np.pi)/2
    ),
    (
        lambda t: 1/2+t/2,
        lambda t: 0
    ),
]
#%%
plt.figure()
plt.grid()
for curve in curves:
    x = [curve[0](t) for t in np.arange(0, 1.01, 0.01)]
    y = [curve[1](t) for t in np.arange(0, 1.01, 0.01)]
    plt.plot(x, y)
#%%
# Точки на границе
border_points = [
    [(curves[0][0](t), curves[0][1](t))
    for t in np.arange(0, 1, 1/(N-1))],
    [(curves[1][0](t), curves[1][1](t))
    for t in np.arange(0, 1, 1/(M-1))],
    [(curves[2][0](t), curves[2][1](t))
    for t in np.arange(0, 1, 1/(M-1))],
    [(curves[3][0](t), curves[3][1](t))
    for t in np.arange(0, 1, 1/(M-1))],
]
# border_points[0]
#%%
# Инициализация сетки
grid = np.zeros((N, M, 2))

# Заполнение граничных условий
for i in range(N-1):
    grid[i, 0, :] = border_points[0][i]
    grid[N-1-i, -1, :] = border_points[2][i]
for j in range(M-1):
    grid[0, M-1-j, :] = border_points[3][j]
    grid[-1, j, :] = border_points[1][j]
#%%
plt.figure()
plt.grid()
for curve in curves:
    x = [curve[0](t) for t in np.arange(0, 1.01, 0.01)]
    y = [curve[1](t) for t in np.arange(0, 1.01, 0.01)]
    plt.plot(x, y)
x = []
y = []
for points in border_points:
    for point in points:
        x.append(point[0])
        y.append(point[1])
plt.plot(x, y, 'o')
plt.plot(grid[:, :, 0], grid[:, :, 1], 'o')
plt.show()
#%%
# Итерационный процесс
for _ in range(100):  # количество итераций
    for i in range(1, N-1):
        for j in range(1, M-1):
            grid[i, j, :] = (k_h*grid[i+1, j, :]
            + k_h*grid[i-1, j, :] + k_v * grid[i, j-1, :]
            + k_v * grid[i, j+1, :]) / (2*k_h+2*k_v)
#%%
plt.figure()
plt.grid()
for i in range(N):
    plt.plot(grid[i, :, 0], grid[i, :, 1], 'o')
for j in range(M):
    plt.plot(grid[:, j, 0], grid[:, j, 1], 'o')
plt.show()
#%%
plt.figure()
plt.grid()
for curve in curves:
    x = [curve[0](t) for t in np.arange(0, 1.01, 0.01)]
    y = [curve[1](t) for t in np.arange(0, 1.01, 0.01)]
    plt.plot(x, y)
x = []
y = []
for points in border_points:
    for point in points:
        x.append(point[0])
        y.append(point[1])
plt.plot(x, y, 'o')
plt.plot(grid[:, :, 0], grid[:, :, 1], 'o')
plt.show()
#%%
plt.figure()
plt.grid()

for i in range(N):
    plt.plot(grid[i, :, 0], grid[i, :, 1])
for j in range(M):
    plt.plot(grid[:, j, 0], grid[:, j, 1])

for i in range(N):
    plt.plot(grid[i, :, 0], grid[i, :, 1], 'o')
for j in range(M):
    plt.plot(grid[:, j, 0], grid[:, j, 1], 'o')

plt.show()
\end{verbatim}

\section{Приложение C}
\begin{verbatim}
import numpy as np
from matplotlib import pyplot as plt
from scipy.optimize import minimize

N = 20
M = 10

# Определение различных фигур
def create_square_curves():
    return [
        (lambda t: 1, lambda t: 2 * t - 1),
        (lambda t: 1 - 2 * t, lambda t: 1),
        (lambda t: -1, lambda t: 1 - 2 * t),
        (lambda t: 2 * t - 1, lambda t: -1)
    ]

def create_arc_curves():
    return [
        (lambda t: np.cos(t * np.pi), lambda t: np.sin(t * np.pi)),
        (lambda t: -1 + t / 2, lambda t: 0),
        (lambda t: np.cos((1 - t) * np.pi)
        / 2, lambda t: np.sin((1 - t) * np.pi) / 2),
        (lambda t: 1 / 2 + t / 2, lambda t: 0),
    ]

def create_circle_curves():
    return [
        (lambda t: np.cos(t * 2 * np.pi), 
        lambda t: np.sin(t * 2 * np.pi)),
        (lambda t: np.cos((t + 1 / 2) * 2 * np.pi),
        lambda t: np.sin((t + 1 / 2) * 2 * np.pi)),
        (lambda t: np.cos((t + 1) * 2 * np.pi), 
        lambda t: np.sin((t + 1) * 2 * np.pi)),
        (lambda t: np.cos((t + 3 / 2) * 2 * np.pi), 
        lambda t: np.sin((t + 3 / 2) * 2 * np.pi)),
    ]

def generate_grid(curves):
    alpha = np.linspace(0, 1, N)
    beta = np.linspace(0, 1, M)
    boundary_points = [
        [(curves[0][0](t), curves[0][1](t)) for t in alpha],
        [(curves[1][0](t), curves[1][1](t)) for t in beta],
        [(curves[2][0](t), curves[2][1](t)) for t in alpha],
        [(curves[3][0](t), curves[3][1](t)) for t in beta],
    ]

    all_points_x = np.zeros((N, M))
    all_points_y = np.zeros((N, M))

    for k in range(N):
        all_points_x[k][0] = boundary_points[0][k][0]
        all_points_y[k][0] = boundary_points[0][k][1]
    for k in range(M):
        all_points_x[N - 1][k] = boundary_points[1][k][0]
        all_points_y[N - 1][k] = boundary_points[1][k][1]
    for k in range(N):
        all_points_x[N - 1 - k][M - 1] = boundary_points[2][k][0]
        all_points_y[N - 1 - k][M - 1] = boundary_points[2][k][1]
    for k in range(M):
        all_points_x[0][M - 1 - k] = boundary_points[3][k][0]
        all_points_y[0][M - 1 - k] = boundary_points[3][k][1]

    all_points_x[1:N-1, 1:M-1] = 
    (all_points_x[1:N-1, 0].reshape(-1, 1) * 
    (M-1 - np.arange(1, M-1)).reshape(1, -1) +                              
    all_points_x[1:N-1, -1].reshape(-1, 1) * np.arange(1, M-1).
    reshape(1, -1)) / (M-1)
    all_points_y[1:N-1, 1:M-1] = 
    (all_points_y[1:N-1, 0].reshape(-1, 1) 
    * (M-1 - np.arange(1, M-1)).reshape(1, -1) +
    all_points_y[1:N-1, -1].reshape(-1, 1) * np.arange(1, M-1).
    reshape(1, -1)) / (M-1)

    return all_points_x, all_points_y

def plot_grid(x, y, title="Grid"):
    plt.figure()
    plt.grid()
    plt.plot(x, y, 'o-')
    plt.plot(x.T, y.T, 'o-')
    plt.title(title)
    plt.show()

# Пример использования
# Квадрат
curves_square = create_square_curves()
grid_x_square, grid_y_square = generate_grid(curves_square)
plot_grid(grid_x_square, grid_y_square,
title="Optimized Grid - Square")

# Арка
curves_arc = create_arc_curves()
grid_x_arc, grid_y_arc = generate_grid(curves_arc)
plot_grid(grid_x_arc, grid_y_arc, title="Optimized Grid - Arc")
\end{verbatim}

\section{Приложение D}
\begin{verbatim}
import numpy as np
from matplotlib import pyplot as plt

def orthogonality_coefficient(grid_x, grid_y):
    N, M = grid_x.shape
    orthogonality = 0
    for i in range(1, N-1):
        for j in range(1, M-1):
            dx1 = grid_x[i+1, j] - grid_x[i, j]
            dy1 = grid_y[i+1, j] - grid_y[i, j]
            dx2 = grid_x[i, j+1] - grid_x[i, j]
            dy2 = grid_y[i, j+1] - grid_y[i, j]
            dot_product = dx1 * dx2 + dy1 * dy2
            norm1 = np.sqrt(dx1**2 + dy1**2)
            norm2 = np.sqrt(dx2**2 + dy2**2)
            angle_cos = dot_product / (norm1 * norm2)
            orthogonality += (1 - angle_cos**2)  # (sin(theta))^2
    return orthogonality / ((N-2) * (M-2))

def uniformity_criterion(grid_x, grid_y):
    N, M = grid_x.shape
    uniformity = 0
    for i in range(1, N):
        for j in range(1, M):
            dx = grid_x[i, j] - grid_x[i-1, j]
            dy = grid_y[i, j] - grid_y[i-1, j]
            distance1 = np.sqrt(dx**2 + dy**2)
            
            dx = grid_x[i, j] - grid_x[i, j-1]
            dy = grid_y[i, j] - grid_y[i, j-1]
            distance2 = np.sqrt(dx**2 + dy**2)
            
            uniformity += abs(distance1 - distance2)
    return uniformity / ((N-1) * (M-1))

def plot_grid(x, y, title="Grid"):
    plt.figure()
    plt.grid()
    plt.plot(x, y, 'o-')
    plt.plot(x.T, y.T, 'o-')
    plt.title(title)
    plt.show()

# Пример использования
N = 20
M = 20

# Определение различных фигур
def create_square_curves():
    return [
        (lambda t: 1, lambda t: 2 * t - 1),
        (lambda t: 1 - 2 * t, lambda t: 1),
        (lambda t: -1, lambda t: 1 - 2 * t),
        (lambda t: 2 * t - 1, lambda t: -1)
    ]

def create_arc_curves():
    return [
        (lambda t: np.cos(t * np.pi), lambda t: np.sin(t * np.pi)),
        (lambda t: -1 + t / 2, lambda t: 0),
        (lambda t: np.cos((1 - t) * np.pi) / 2,
        lambda t: np.sin((1 - t) * np.pi) / 2),
        (lambda t: 1 / 2 + t / 2, lambda t: 0),
    ]

def create_circle_curves():
    return [
        (lambda t: np.cos(t * 2 * np.pi),
        lambda t: np.sin(t * 2 * np.pi)),
        (lambda t: np.cos((t + 1 / 2) * 2 * np.pi),
        lambda t: np.sin((t + 1 / 2) * 2 * np.pi)),
        (lambda t: np.cos((t + 1) * 2 * np.pi),
        lambda t: np.sin((t + 1) * 2 * np.pi)),
        (lambda t: np.cos((t + 3 / 2) * 2 * np.pi),
        lambda t: np.sin((t + 3 / 2) * 2 * np.pi)),
    ]

def generate_grid(curves):
    alpha = np.linspace(0, 1, N)
    beta = np.linspace(0, 1, M)
    boundary_points = [
        [(curves[0][0](t), curves[0][1](t)) for t in alpha],
        [(curves[1][0](t), curves[1][1](t)) for t in beta],
        [(curves[2][0](t), curves[2][1](t)) for t in alpha],
        [(curves[3][0](t), curves[3][1](t)) for t in beta],
    ]

    all_points_x = np.zeros((N, M))
    all_points_y = np.zeros((N, M))

    for k in range(N):
        all_points_x[k][0] = boundary_points[0][k][0]
        all_points_y[k][0] = boundary_points[0][k][1]
    for k in range(M):
        all_points_x[N - 1][k] = boundary_points[1][k][0]
        all_points_y[N - 1][k] = boundary_points[1][k][1]
    for k in range(N):
        all_points_x[N - 1 - k][M - 1] = boundary_points[2][k][0]
        all_points_y[N - 1 - k][M - 1] = boundary_points[2][k][1]
    for k in range(M):
        all_points_x[0][M - 1 - k] = boundary_points[3][k][0]
        all_points_y[0][M - 1 - k] = boundary_points[3][k][1]

    all_points_x[1:N-1, 1:M-1] = (all_points_x[1:N-1, 0]
    .reshape(-1, 1) * (M-1 - np.arange(1, M-1)).reshape(1, -1) 
    all_points_x[1:N-1, -1].reshape(-1, 1)
    * np.arange(1, M-1).reshape(1, -1)) / (M-1)
    all_points_y[1:N-1, 1:M-1] = (all_points_y[1:N-1, 0]
    .reshape(-1, 1) * (M-1 - np.arange(1, M-1)).
    reshape(1, -1) + all_points_y[1:N-1, -1].
    reshape(-1, 1) * np.arange(1, M-1).
    reshape(1, -1)) / (M-1)

    return all_points_x, all_points_y

# Пример использования
# Квадрат
curves_square = create_square_curves()
grid_x_square, grid_y_square = generate_grid(curves_square)
plot_grid(grid_x_square, grid_y_square,
title="Optimized Grid - Square")

# Арка
curves_arc = create_arc_curves()
grid_x_arc, grid_y_arc = generate_grid(curves_arc)
plot_grid(grid_x_arc, grid_y_arc, title="Optimized Grid - Arc")

# Круг
curves_circle = create_circle_curves()
grid_x_circle, grid_y_circle = generate_grid(curves_circle)
plot_grid(grid_x_circle, grid_y_circle,
title="Optimized Grid - Circle")

# Вычисление коэффициентов
orthogonality_square = orthogonality_coefficient
(grid_x_square, grid_y_square)
uniformity_square = uniformity_criterion
(grid_x_square, grid_y_square)

orthogonality_arc = orthogonality_coefficient
(grid_x_arc, grid_y_arc)
uniformity_arc = uniformity_criterion
(grid_x_arc, grid_y_arc)

orthogonality_circle = orthogonality_coefficient
(grid_x_circle, grid_y_circle)
uniformity_circle = uniformity_criterion
(grid_x_circle, grid_y_circle)

print(f"Square Grid - Orthogonality: {orthogonality_square}, 
Uniformity: {uniformity_square}")
print(f"Arc Grid - Orthogonality: {orthogonality_arc}, 
Uniformity: {uniformity_arc}")
print(f"Circle Grid - Orthogonality: {orthogonality_circle},
Uniformity: {uniformity_circle}")
\end{verbatim}
\end{document}
